% =====================================================================
%  Capitolo 5 — Dalla ricerca alle idee: competitor, perceptual mapping,
%  analisi di utenti e mercato, motivazioni dei giocatori e brainstorming
%  (problema e soluzione, seconda parte)
%  Fonte: trascrizione della lezione GD5 (Game Design, primo anno)
%  Pacchetti richiesti nel preambolo del libro:
%    \usepackage[italian]{babel}, \usepackage[T1]{fontenc},
%    \usepackage[utf8]{inputenc} (solo pdfLaTeX), \usepackage{booktabs}
%  Nota sulla fonte: la trascrizione automatica storpia diversi nomi
%  propri (titoli, studi, autori) e contiene alcuni tratti ripetuti in
%  loop o incomprensibili (durante la pausa e in un paio di battute in
%  aula). Dove il riferimento è ricostruibile con certezza lo si riporta
%  nella forma corretta; dove non lo è, lo si segnala in nota invece di
%  indovinare.
% =====================================================================

\chapter[Dalla ricerca alle idee]{Dalla ricerca alle idee: perceptual
mapping, utenti, mercato, motivazioni e brainstorming}
\label{cap:ricerca-idee}

Il capitolo precedente si è fermato a metà strada: la lezione su
\emph{problema e soluzione} era stata annunciata come troppo lunga per una
sola giornata, e si era chiusa con le domande per l'analisi dei competitor
(sezione~\ref{sec:analisi-competitor}). Questo quinto capitolo la porta a
termine. L'obiettivo dichiarato dal docente è ``realisticamente'' solo
questo, finire la lezione; e in effetti la giornata copre tutto il resto
della ricerca in pre-produzione, fino alla sua conclusione naturale (la
\emph{visione}), e apre la parte successiva, quella sulla
\textbf{generazione delle idee}.

La lezione è tenuta ancora da Mattia Podini. Il percorso del capitolo ne
ricalca la struttura:
\begin{enumerate}
  \item le domande in apertura: il \textbf{production dictionary} ripreso e
        approfondito, con una digressione su come è organizzata la
        documentazione nelle produzioni moderne e sulla distinzione tra
        documenti di visione e documenti di produzione
        (sezioni~\ref{sec:pd-ripreso} e~\ref{sec:documenti-visione});
  \item una breve ripresa del punto in cui ci si era fermati, con una
        parentesi su Deming (sezione~\ref{sec:ripresa-gd5});
  \item gli ultimi strumenti per l'analisi dei competitor: i framework usati
        anche per l'analisi, l'\textbf{elicitazione triadica} e la
        \textbf{perceptual mapping}, con due esercizi in aula
        (sezioni~\ref{sec:framework-triadica} e~\ref{sec:perceptual-mapping});
  \item le altre fonti di dati, l'\textbf{analisi degli utenti} e
        l'\textbf{analisi di mercato}
        (sezioni~\ref{sec:fonti-dati}, \ref{sec:analisi-utenti}
        e~\ref{sec:analisi-mercato});
  \item l'analisi di gameplay attraverso il modello delle
        \textbf{motivazioni} di Quantic Foundry
        (sezione~\ref{sec:quantic-foundry});
  \item la \textbf{correlazione} dei dati e la chiusura della pre-produzione
        con le domande che portano alla \emph{game vision}
        (sezioni~\ref{sec:correlazione} e~\ref{sec:fine-preproduzione});
  \item l'inizio della parte su \textbf{idee e brainstorming}: i principi, il
        \emph{design elegante}, la creatività e una rassegna di tecniche
        (sezioni~\ref{sec:principi-brainstorming}--\ref{sec:tecniche-idee});
  \item il task della settimana, a gruppi (sezione~\ref{sec:task-gd5}).
\end{enumerate}

Come per il capitolo precedente, la lezione non arriva fino in fondo: la
parte su \emph{come} si conduce un brainstorming è rimandata alla lezione
del venerdì successivo.

% ---------------------------------------------------------------------
\section{In apertura: il production dictionary, ripreso}
\label{sec:pd-ripreso}

\subsection{Prima di cominciare}

La registrazione si apre sulla coda di un momento informale: il docente sta
commentando una partita, e osserva che una \emph{Blue Map} ``non sarebbe per
niente male'', perché permetterebbe di vedere subito dove sono le stanze
segrete senza sprecare bombe per cercarle; resta da trovare ``al più
presto un'altra monetina'' che gliela consigli.\footnote{Il riferimento è
quasi certamente a \emph{The Binding of Isaac}, in cui la \emph{Blue Map}
è un oggetto che rivela le stanze segrete della mappa. Il contesto (una
partita in corso prima della lezione, forse con suggerimenti dalla classe)
non è ricostruibile dalla trascrizione.}

Poi si entra nel vivo, con il solito rituale: prima di cominciare, domande o
dubbi sulla lezione precedente, sui task in generale, su qualsiasi cosa. Il
docente ricorda che il termine ultimo per la consegna dei \textbf{task
sugli arcade} è scaduto: i documenti restano lì, ``fermi, fissi'', finché non
li commenterà, probabilmente nel fine settimana (si veda il task in
sezione~\ref{sec:task-gd2}).

\paragraph{Un altro brevetto aggirato.} Prima della lezione alcuni studenti
stavano discutendo ``con grande fervore'' di un altro esempio di brevetto
aggirato, sulla scia del caso Pokémon visto la volta precedente
(sezione~\ref{sec:analisi-competitor}). Il docente conferma che di esempi
simili ce n'è più di uno, ma che non serve approfondirli; non sa che cosa sia
stato detto prima del suo arrivo.

\subsection{Nessun limite di pagine}

Uno studente chiede di ripetere in cinque minuti il concetto di
\textbf{production dictionary} (sezione~\ref{sec:production-dictionary}).
``Certamente.''

Lo scopo di un production dictionary è lo stesso per cui esiste un
dizionario: \textbf{spiegare terminologie specifiche}. Nelle produzioni vere
non esiste una limitazione come quella data nel task: non esiste un
production dictionary ``di una pagina'', ``di due pagine'' o comunque
limitato in grandezza. Il limite di una pagina era un vincolo
dell'esercizio (sezione~\ref{sec:task-gd4}).

\subsection{File a sé o paragrafo d'apertura}

La discriminante più grossa, come detto l'ultima volta, è \textbf{dove sta}
il production dictionary nella struttura della documentazione. Il docente
nota che è un argomento che non si tratta quasi mai, se lo segna e promette
di aggiungere qualche slide.
\begin{description}
  \item[Come documento a sé] può restare un file autonomo all'interno del
        Drive (``assumiamo che la vostra azienda lavori con Google Drive,
        cosa possibilmente non vera'').
  \item[Come elemento di un altro documento] può diventare un paragrafo, o
        un'intestazione, ``non so come volete chiamarla'', posto
        \textbf{in apertura} di un singolo documento, per aiutare il lettore
        a capire in anticipo eventuali termini.
\end{description}

\paragraph{L'esempio di Wikipedia.} Wikipedia funziona più o meno così. Gli
argomenti sono larghissimi e diversi, ma ci sono sezioni che indirizzano a
voci specifiche. Si prenda la pagina di \emph{Liverpool}: tutte le
informazioni al suo interno, o già la colonna a sinistra, rimandano a cose
che si dovrebbero sapere per la buona comprensione di quella pagina, come
``che cos'è una città'' o ``che cos'è una bandiera''. Sono cose estremamente
pratiche e semplici, ma il principio è lo stesso.

\paragraph{L'esempio del GDD dello sparo.} Si supponga di scrivere un GDD
(un \emph{design document}: ``GDD'' qui vale per qualsiasi disciplina, e
diventa il design document di quella disciplina) dedicato a una meccanica,
per esempio il salto. In un gioco come \emph{Mario} si può saltare e sparare:
è possibile, non necessario, che nel documento si voglia spiegare anche la
meccanica di sparare (``datemi un salto, lo sparo'', si corregge il docente
dopo essersi confuso tra le due). La si può spiegare nel production
dictionary in testa al documento oppure, in alternativa, \textbf{linkare} e
rimandare al documento ``GDD Sparo''. Il production dictionary aiuta proprio
in questo: se ci sono terminologie necessarie alla buona comprensione di
\emph{quel} documento, le si chiarisce lì.

\paragraph{Attenzione al lettore.} Resta valido quanto detto l'altra volta:
un \textbf{blocco di terminologia} all'inizio di un documento probabilmente
non è piacevole da vedere per il lettore (sezione~\ref{sec:golden-rules}).
Le soluzioni sono quindi due: un documento a sé, che sta in un'altra parte
rispetto agli altri, oppure un production dictionary \textbf{ibridato} con
altre cose.

\subsection{Ibridare i documenti}

Il production dictionary non è l'unico caso: molti altri documenti, anzi in
realtà tutti, si possono ibridare. Nessuno vieta di prendere i punti che un
documento dovrebbe svolgere o contenere (le categorie che si vedranno per gli
altri tipi di documento) e mescolarli con altro: ``si fa''. Il documento che
più si presta all'ibridazione è proprio il \textbf{GDD}, perché non ha una
forma specifica: è malleabile (si veda anche quanto detto sui template,
sezione~\ref{sec:due-principi}).

\subsection{Che cosa ci va, e chi lo decide}

Il production dictionary deve spiegare i \textbf{termini necessari per il
prodotto}, termini che tutte le persone che ci lavorano devono condividere.
Il \emph{mana}, se il gioco lo ha, può sembrare un termine che non serve
spiegare; ma il docente assicura che, parlando con un sound designer o con
figure più lontane dal design, le cose cambiano. Molti collaboratori sono
\emph{outsourced}, presi al di fuori del team; lui stesso ha collaborato con
tanti animatori che avevano poca conoscenza videoludica, perché venivano dal
cinema o dall'animazione e non erano designer. Si assume che, bene o male,
tutti giochino ai videogiochi (vale anche per gli artisti), ma spesso non è
detto. Alcune terminologie che per un designer sono comuni devono esserlo
\textbf{per tutti}, e quindi vanno spiegate.

\paragraph{Chi sceglie i termini.} Che cosa mettere nel production dictionary
lo sceglie il designer. In un documento specifico come quello del task non lo
sceglie il junior: è il \textbf{lead} a decidere quali termini inserire. A
riempirlo, però, sono poi molte mani diverse: chi ha fatto il documento del
salto aggiunge le sue terminologie specifiche, chi ha fatto quello dello
sparo avrà altre idee e le inserirà a sua volta, qualora servano.

\paragraph{Uno per progetto.} Come scritto anche nelle slide, questi documenti
possono avere un corrispettivo in altri dipartimenti (il \emph{lexicon},
sezione~\ref{sec:production-dictionary}). Ma c'è un'informazione che il
docente non aveva scritto, e che dovrà aggiungere: di production dictionary
ce n'è \textbf{uno per progetto}. È uno di quei documenti di cui ne esiste
uno solo, e che viene sempre aggiornato.

% ---------------------------------------------------------------------
\section{Documenti di visione e documenti di produzione}
\label{sec:documenti-visione}

\subsection{Una categoria personale}

Quanto appena detto vale, in generale, per tutti i documenti che il docente
chiama \textbf{documenti di visione}. Avverte che non è una categoria
largamente condivisa: alcuni documenti ci rientrano, altri no, ma nella sua
testa ``aiuta, mi dà ordine mentale'', e potrebbe essere utile anche agli
studenti.
\begin{description}
  \item[Documenti di visione] Rimangono ``nei secoli dei secoli'': vengono
        aggiornati, certo, ma aiutano a capire il progetto ad
        \textbf{altissimo livello}. Ne fanno parte il production dictionary
        (per la comprensione delle nomenclature), il \textbf{documento di
        visione} già accennato più volte, che si vedrà (uno per progetto:
        quando viene creato, ``quello è e quello rimane''), ma anche i
        \textbf{pitch}, di cui esisterà magari qualche versione. Sono
        documenti che servono a vendere il progetto, o che ne incorporano
        l'idea, il concept, il cuore.
  \item[Documenti di produzione] Sono i documenti che servono effettivamente
        a produrre il gioco: i \textbf{GDD} sono documenti di produzione.
        Ce ne sono anche altri che, per esempio, tracciano lo stato della
        produzione.
\end{description}

\subsection{Dalla Bibbia alla wiki, e oltre}

Il punto che chiarisce molte cose, e che l'ultima volta non era stato
trattato, è che nelle produzioni moderne i documenti sono
\textbf{separati}. Il grosso GDD, l'unico documento che spiega tutto il gioco
\emph{in toto}, a livello produttivo \textbf{non esiste più} (si veda
l'approfondimento sul mito del Game Design Document,
sezione~\ref{sec:due-principi}). Sopravvive solo nel reame del marketing, o
in approcci che non fanno parte della produzione e che qui non interessano.

Lo step intermedio sono state le \textbf{wiki di design}: è il caso di chi usa
\emph{Confluence}, che ``è proprio questo''. Dalla ``Bibbia'' si salva quindi
la wiki, e poi c'è ciò che si usa oggi, che il docente chiama, in mancanza di
un nome migliore, ``una specie di wiki 2''. Dipende dalle aziende; ma
lavorando dentro Google Drive, gli studenti fanno già qualcosa di simile.

\paragraph{Documenti che smistano i lettori.} In questa struttura c'è una
documentazione di visione, di organizzazione del progetto (o di
pre-produzione, ``come vi pare'') che guida chi lavora al progetto verso gli
altri documenti. Aiuta la comprensione del progetto oppure reindirizza il
lettore verso le parti specifiche che gli servono:
\begin{itemize}
  \item un \textbf{artista} viene spedito nelle cartelle in cui sono
        archiviati gli asset;
  \item un \textbf{programmatore} viene indirizzato, per esempio, alla
        cartella in cui viene messa la build (se la si carica sul Drive);
  \item chi deve \textbf{implementare una meccanica}, programmatore o artista
        che sia, viene rimandato al documento specifico: deve implementare il
        salto, va al documento del salto; deve implementare lo sparo, va a
        quello dello sparo.
\end{itemize}

\subsection{Chi tiene in ordine il ``giardinetto''}

La gerarchia di base è molto semplice; la sotto-gerarchia della
sotto-gerarchia spetta a chi gestisce lo spazio di archiviazione: il docente
dice ``Drive'', ma può essere qualsiasi storage di dati o file, ``potrebbe
essere pure su Git''. Nel corso se ne occupano gli studenti, perché sono
designer e il cerchio si restringe; ma in una grande produzione è
responsabilità della \textbf{produzione}, cioè del team di produzione,
generalmente composto da producer e production director (dipende da quanti
sono). Sono loro a strutturare le cartelle e l'organizzazione del Drive, di
Confluence o di qualsiasi altro strumento usato per raccogliere la
documentazione. Poi ognuno lavora \textbf{nel proprio giardinetto}, che
naturalmente va mantenuto ordinato.

È esattamente ciò che succede nel corso: ci sono delle regole, i documenti si
creano nelle cartelle dei task (GDLD e così via), ed è il docente a spostarli
quando serve (sezione~\ref{sec:drive-regole}). ``Il rapporto è lo stesso,
uno a uno.''

Il docente nota il silenzio della classe (``non so se è un bene o un male;
assumiamo che sia un bene''), invita a parlare in caso di perplessità e si
torna alla lezione.

% ---------------------------------------------------------------------
\section{Dove eravamo rimasti}
\label{sec:ripresa-gd5}

\subsection{Il mantra del compromesso}

La volta precedente si era iniziato il discorso su come individuare il
problema o la necessità e come categorizzarli. Prima di proseguire, il
docente torna sulla slide ``più importante'', il \textbf{mantra} entro il
quale la produzione agisce: tutti gli ingranaggi si muovono e il
\textbf{compromesso} determina che cosa si dovrebbe fare a livello di design,
qual è la miglior soluzione possibile (sezione~\ref{sec:design-compromesso}).
Si erano già viste la citazione di Diego, il direttore del corso, e quella di
Deming (sezione~\ref{sec:dati}).

\subsection{Una parentesi su Deming}

Per chi volesse più informazioni, il docente lascia due ``ganci''. Non si può
parlare di tutto (``mi piacerebbe, ma non si può''), quindi questi servono a
fare ricerca per conto proprio.

\paragraph{Il ciclo di Deming.} La metodologia di Deming è un modo di fare
produzione basato su un ciclo: si pianifica, si fa qualcosa, si verifica, si
agisce, e si ricomincia.\footnote{In aula l'ordine delle fasi viene detto in
modo confuso (``check, act, plan, do''). La forma canonica del ciclo di
Deming, o ciclo PDCA, è \emph{Plan--Do--Check--Act}: pianificare, eseguire,
verificare, agire (standardizzando ciò che funziona e correggendo ciò che non
funziona), e poi ripartire.} Il docente scherza anche sulla foto di Deming
nelle slide (``bellissimo uomo'').

\paragraph{Le sette malattie mortali.} Un'altra cosa interessante di Deming
sono le \textbf{sette malattie mortali della gestione occidentale}. Deming
era americano, e poi si è spostato in Oriente (in Giappone); alcune di queste
malattie risentono della sua cultura: al punto sei, per esempio, ci sono i
\emph{costi medici eccessivi}. Altre, invece, potrebbero essere rilevanti da
approfondire.\footnote{L'elenco, come formulato da Deming in \emph{Out of the
Crisis} (1982), è: (1) mancanza di costanza nello scopo; (2) enfasi sui
profitti a breve termine; (3) valutazione delle prestazioni, \emph{merit
rating} o revisione annuale; (4) mobilità del management; (5) dirigere
un'azienda basandosi solo sui dati visibili; (6) costi medici eccessivi;
(7) costi eccessivi di garanzia e contenziosi legali.} Quella su cui si
sofferma il docente, legata anche al terzo punto, è \textbf{dirigere
un'azienda basandosi solo sui dati visibili e misurabili}: il fatto che la
produzione possa non valutare il benessere o il morale del singolo individuo,
degli operativi. Il docente lo ha visto spesso, ed è effettivamente un
problema nelle produzioni. È uno dei motivi per cui molte produzioni
\emph{indie} sono più sostenibili: essendo in numero minore, ci si può
preoccupare umanamente delle persone con cui si lavora. ``Tuttavia, non è
sempre vero.''

\subsection{La mappa della ricerca, a che punto siamo}

Si torna alla \textbf{ricerca} e all'esempio guida della lezione precedente:
il ``magico'' Poké Walker modificato, mettiamola così, per curare l'obesità
infantile (sezione~\ref{sec:dati}). Degli otto punti della mappa
(sezione~\ref{sec:mappa-ricerca}) si erano visti:
\begin{itemize}
  \item i competitor diretti, indiretti e similari
        (sezione~\ref{sec:competitor-diretti}); altri strumenti si vedranno
        anche nella parte di brainstorming e di generazione delle idee;
  \item la definizione degli stakeholder (sezione~\ref{sec:stakeholder});
  \item la definizione di problema contro necessità (``o comunque, se c'è un
        impiccio: non ho un sinonimo migliore''), con l'identificazione della
        categoria entro cui la questione si presenta, che dovrebbe aiutare a
        capire come affrontarla: nella maniera più efficiente e veloce
        possibile, o in una maniera più costosa ma potenzialmente con grandi
        ritorni a lungo termine (sezione~\ref{sec:problema-necessita});
  \item altre metodologie e categorizzazioni, e infine le domande
        sull'analisi dei competitor (sezione~\ref{sec:analisi-competitor}).
\end{itemize}
Come al solito, avverte il docente, in questa lezione soprattutto arriva una
``valanga'' di strumenti potenzialmente utili. Non vanno usati tutti: è
irrealistico pensare di usarli sempre tutti.

% ---------------------------------------------------------------------
\section{Analisi dei competitor: framework ed elicitazione triadica}
\label{sec:framework-triadica}

\subsection{I framework servono anche all'analisi}

Le domande viste l'ultima volta sono domande \textbf{ad alto livello} da cui
partire per analizzare i competitor, o, più realisticamente, alcuni aspetti
o sistemi di altri giochi. Per farlo si usano \textbf{framework adatti}: i
framework visti in design, ma non solo (anche quelli di level design), si
possono applicare sia all'analisi sia al design, e spesso ad entrambe
contemporaneamente. Gli studenti hanno appena usato i \emph{core elements}
per analizzare gli arcade; li possono usare anche per progettare
(sezione~\ref{sec:romero-schreiber}).

\paragraph{Anche per il pitch.} Quando il docente fa una presentazione, per
il proprio gioco o per quello di qualcun altro (gli è capitato di fare
\emph{pitching} anche per altri: esiste anche questa figura), una delle
prime slide serve a far capire il gioco all'ascoltatore. Come? Per esempio
con una lista di \textbf{game qualities} di Romero e Schreiber: che tipo di
camera ha, che genere è, se è multiplayer o single player, eventuali altre
cose. Lo stesso strumento serve quindi all'analisi, al pitching e al design.
I \emph{core elements} sono un framework ancora molto basilare, così come
quello di Bura (sezione~\ref{sec:bura}); il discorso su come fare design con
quest'ultimo era già stato fatto, è complesso, e il docente invita prima a
fare ricerca: sul sito di Bura si trova tutto.

\subsection{L'elicitazione triadica}

Un altro modo per analizzare i competitor ha un nome ``estremamente
complesso, per dire una cosa abbastanza semplice'': l'\textbf{elicitazione
triadica} (\emph{triadic elicitation}).\footnote{Nella trascrizione il termine
compare storpiato (``licizzazione triadica''). L'elicitazione triadica è una
tecnica di intervista nata in psicologia (nell'ambito della \emph{repertory
grid} di George Kelly) e usata anche nella ricerca di mercato e nella user
research.}
\begin{enumerate}
  \item Si individuano i \textbf{tre competitor principali}: solo tre, sempre
        per una questione di compromesso, cioè del tempo da spendere in
        ricerca, che spesso non basta.
  \item Si chiede al target audience di \textbf{raggruppare i due più
        simili}.
  \item Si chiede \textbf{perché} sono simili e in che cosa sono diversi;
        che cosa piace di quei due, che cosa non piace.
  \item Del gioco rimasto fuori, quello differente, si chiede perché è
        diverso dagli altri, in che cosa è simile, se è meglio o peggio (in
        alcuni aspetti sarà meglio, in altri peggio) e perché.
\end{enumerate}

\paragraph{La ``piaga'' dei soulslike.} C'è sempre l'esempio dei soulslike
(``non che sia problematico, ma è sempre quello''): che cosa rende questi
prodotti dei soulslike? Che cosa li fa rientrare in quel genere? Oppure: qual
è quel \emph{core element} che si sta cercando, quell'elemento che poi
determinerà anche il proprio prodotto? L'elemento importante è la
\textbf{precisione}: bisogna cercare di prendere l'elemento veramente
\textbf{dominante}.

\subsection{Parlare con l'utenza}

È un primo punto: \textbf{parlare con l'utente}. Non ci si va senza nulla: ci
si premunisce, si procurano prima i dati ed eventualmente già una lista di
competitor. Poi si chiede al target audience quali sono i \textbf{pain
point}, cioè le problematiche che individua, e anche che cosa gli piace,
per arrivare potenzialmente all'\textbf{insight}: la consapevolezza su ciò
che si è ottenuto.

\paragraph{Il rischio di allontanarsi dall'utente.} Il metodo aiuta anche a
capire come i giochi vengono \textbf{percepiti dall'utenza}. Più si fa design,
più tempo si dedica a progettare giochi, più c'è il rischio di
``sfilarsi'' dal ruolo di utente e diventare soltanto creatori di contenuti;
a meno di non fare proprio il gioco a cui si gioca ventiquattro ore su
ventiquattro. Se il docente dovesse fare un gioco di guida, sarebbe molto
lontano da quell'audience: analizzando i giochi da designer, potrebbe leggere
una meccanica in un modo che non corrisponde affatto a come la percepisce
l'audience. L'unico modo è, ancora una volta, mettersi nei panni
dell'audience; e, come detto l'ultima volta, è difficile
(sezione~\ref{sec:utenti-mercato}).

% ---------------------------------------------------------------------
\section{La perceptual mapping}
\label{sec:perceptual-mapping}

\subsection{Che cos'è}

Un altro metodo, molto utile, è ``il cavallo di battaglia'' di tanti
designer e di tante messe a terra di idee: funziona molto bene sia nel design
sia nella spiegazione. È la \textbf{perceptual mapping} (il docente non
ricorda la traduzione italiana esatta).\footnote{In italiano si parla di
\emph{mappa percettiva} o \emph{mappa di posizionamento}.}

Si prendono anche qui le \textbf{game qualities} principali dei competitor, o
di un mercato in generale su cui ci si sta affacciando, e le si dispone su
un grafico a due assi per individuare i \textbf{vuoti}: i punti in cui non
c'è nessuno, cioè, fondamentalmente, i \textbf{vuoti di mercato}.

\paragraph{Gli assi sono spettri.} Ogni asse deve essere uno
\textbf{spettro}: va da un estremo, il più possibile verso 1, all'altro, il
più possibile verso 0. Le slide mostrano due esempi lontani dai videogiochi.
\begin{description}
  \item[Automobili] Sugli assi ci sono ``conservativo'' contro ``sportivo''
        (che, a rigore, non sono proprio uno il contrario dell'altro) e
        ``vistoso, distintivo'' contro ``pratico, economico'', che invece
        stanno bene o male ai due estremi.
  \item[Farmaci] Sugli assi ci sono l'efficacia e l'impatto sul corpo. Un
        antibiotico ha un effetto molto importante sul sistema immunitario;
        un fluidificante (``abbiamo fatto pubblicità occulta'') è meno
        impattante; le caramelle alla menta, le Golia, ancora meno. Se
        l'obiettivo fosse un farmaco estremamente efficace che però
        ``distrugga completamente il corpo del paziente'', quel quadrante è
        coperto.
\end{description}

\paragraph{Un vuoto non è un invito.} Nell'esempio dei farmaci c'è un intero
quadrante vuoto: poco efficace e molto ``gentile''. Probabilmente c'è un
motivo se è vuoto. Aver identificato un vuoto non significa doverci
``fiondarsi a capofitto'': è probabile che non ci sia grande richiesta per
qualcosa di molto delicato e poco efficace. Qui alcuni studenti obiettano che
in realtà quel mercato esiste, e il docente concede (``non in contesti
farmaceutici specifici, però esiste, sì, avete ragione''): è un altro
mercato. Si tornerà sul tema con la correlazione dei dati
(sezione~\ref{sec:correlazione}).

\paragraph{I due obiettivi.} Lo stesso lavoro si fa sui giochi, e gli
obiettivi principali per cui si fa una perceptual mapping sono due:
\begin{itemize}
  \item trovare \textbf{combinazioni innovative};
  \item trovare \textbf{buchi di mercato}.
\end{itemize}

\subsection{La perceptual mapping nel design: i nemici}

Il docente la usa soprattutto non per fare analisi, ma per fare
\textbf{design}. Non crede che nel corso ci sarà modo di parlare di
\emph{enemy design} (forse nel \emph{core combat}), ma cita \emph{Doom} come
caso di un roster di nemici costruito proprio così.\footnote{Il passaggio
nella trascrizione è disturbato (``se non è a Ro Doom, è stato fatto proprio
così: i nemici di Doom sono stati fatti così''). Si riporta il riferimento a
\emph{Doom} senza poter precisare quale capitolo della serie intendesse il
docente.} L'anno successivo, forse, ci sarà l'\emph{encounter design} con
Stefano; ma l'encounter design si basa sul \textbf{posizionamento} dei
nemici, più che sulla loro tipologia, anche se il comportamento viene in
parte toccato. ``Apro e chiudo questa parentesi.''

Il docente usa il metodo per trovare \textbf{combinazioni di nemici} che non
sono ancora state fatte, o qualcosa che \textbf{spinga il giocatore} nella
direzione voluta, lato gameplay. L'esempio:
\begin{itemize}
  \item su un asse, \textbf{DPS} contro \textbf{supporto};
  \item sull'altro, tanti HP (il \textbf{tank}) contro pochi HP: non
        necessariamente \emph{glass cannon}, suggerisce uno studente, ma
        \emph{squishy}.
\end{itemize}

\paragraph{Il glass cannon.} Un nemico con tanti danni e pochi HP è un
\emph{glass cannon}, un ``cannone di vetro'': strategie, personaggi o nemici
che fanno tantissimo danno ma che, se li si tocca, ``saltano proprio per
aria''. Uno studente propone un esempio da una build con un \emph{charm} (il
riferimento sembra a \emph{Hollow Knight}); il docente concede che non è
proprio un glass cannon, perché ``chiaramente non si può colpire'' (``siete
dei disonesti''). Esempi più netti sono i \textbf{cecchini}, anche in
\emph{XCOM}: tutti quei nemici che fanno tanti danni e la cui strategia,
generalmente, è difendersi dietro ai tank, che hanno tanta vita ma fanno meno
danni: assorbono, ma non colpiscono duro.

\paragraph{I punti ciechi.} Mappati i nemici, si vedono i \textbf{punti
ciechi}: magari manca un nemico di supporto con tanti punti vita. Magari non
interessa, magari sì. Generalmente, soprattutto nell'enemy design,
interessano gli \textbf{estremi}, perché sono quelli che \textbf{spingono il
giocatore}.

\paragraph{Un grafico non basta.} È l'altro punto ``importantissimo'', da
imparare a memoria: \textbf{un grafico cartesiano con due assi non basta, ne
servono più di uno}. Quello appena fatto non mappa tutti i nemici: il glass
cannon potrebbe essere lento o veloce, e allora serve un asse \emph{slow} --
\emph{fast}. In aula se ne aggiungono altri:
\begin{itemize}
  \item la \textbf{grandezza}: interessa relativamente, può essere utile per
        comunicare al giocatore ma è più un discorso artistico; a meno che il
        gioco non abbia una meccanica di \emph{shrinking} o non si basi
        proprio sulle dimensioni, e allora sì;
  \item \textbf{melee} contro \textbf{ranged} (``bravo'').
\end{itemize}
Il cecchino, glass cannon, sarà ranged. A rigore di logica andrebbe messo su
\emph{slow}; ma se ci fosse un altro nemico con tanti DPS e tanti HP (che
quindi non è più un glass cannon), probabilmente sarebbe quello ad andare su
\emph{slow}, e il cecchino su \emph{fast}, o a metà: andrebbe capito insieme
agli altri nemici.

\paragraph{Combinare le categorie.} Già così si possono costruire categorie di
nemici che, combinate, danno \textbf{sfide diverse} al giocatore:
\begin{itemize}
  \item il cecchino che si mette dietro il tank: la sfida per il giocatore è
        \textbf{aggirare il cecchino} (il docente lo ripete perché Mattia,
        collegato da remoto, possa seguire meglio);
  \item il supporto, per cui la dinamica di gioco è bene o male sempre la
        stessa;
  \item il classico nemico che fa un sacco di danni e ti corre addosso per
        farsi saltare in aria, di solito anche lui \emph{squishy};
  \item il tank con tanta vita che fa poco danno; oppure il tank con tanta
        vita, poco danno e molto lento; e così via.
\end{itemize}

\subsection{Esercizio in aula: la mappa dei metroidvania}

Tornando all'analisi, si prova a farne una con la classe (``tanto il tempo lo
abbiamo''). Si vuole fare un ``nuovo, incredibile, innovativo gioco''.
Naturalmente non si parte dal nulla: si restringe il campo, a partire da un
\textbf{genere}. Il genere scelto è il \textbf{metroidvania}.

\paragraph{Gli assi.} Il docente chiede alla classe che cosa mettere sugli
assi (``probabilmente qualsiasi cosa, veramente''). Non deve per forza esserci
correlazione tra un asse e l'altro; ciò che conta è che ciò che si include
sia \textbf{in comune} a tutti i giochi mappati. Emergono:
\begin{itemize}
  \item \textbf{narrative first} contro \textbf{gameplay first} (proposto
        come ``spiccatamente narrativo'');
  \item \textbf{high challenge} contro \textbf{low challenge}, o
        accessibilità;
  \item \textbf{prima persona} contro \textbf{terza persona}, e \textbf{2D}
        contro \textbf{3D}: un'ottima scelta se si dovesse fare ricerca,
        perché dà già un'idea, oltre che del discorso di gameplay, anche di
        quello artistico;
  \item \textbf{level design focused} contro \textbf{combat focused} (``LD
        focused, combat focused'').
\end{itemize}
Il docente suggerisce altri assi possibili, che però si escludono dal grafico
per non complicarlo troppo: la durata del gameplay; la \textbf{progressione},
che può essere basata sul \emph{grinding} o no, e che può essere una
progressione di \emph{efficienza} (il giocatore ha sempre le stesse meccaniche
ma deve diventare più efficiente nell'usarle) o di \emph{potere} (\emph{power
progression}, ci si tornerà con le motivazioni,
sezione~\ref{sec:quantic-foundry}). Anche la camera e il controllo
emergono, ma ``inizia a diventare un po' complesso''.

\paragraph{Una convenzione percettiva.} Nel disegnare il grafico il docente
nota di aver ``sbagliato'': c'è una convenzione percettiva per cui il valore
\emph{alto} si mette in su o a destra. ``Ce ne faremo una ragione'': nessuno
vieta di \textbf{numerare} gli assi, di dare delle tacche. I valori vanno
trovati in modo approssimativo, ma sempre con un discorso di precisione.

\paragraph{I giochi.} Si comincia a riempire, ``a spanne'' (``chiaramente ci
servirebbero dei dati'').
\begin{itemize}
  \item I primi tre \emph{Castlevania} vengono proposti, ma il docente
        obietta che non sono metroidvania nel senso pieno.\footnote{La
        trascrizione è qui disturbata (``le prime tre Castlevania non sono
        mica in bagno''); il senso più probabile è che i primi tre
        \emph{Castlevania} non siano ancora metroidvania, genere che nasce
        con \emph{Castlevania: Symphony of the Night} e con la serie
        \emph{Metroid}.}
  \item \emph{Hollow Knight}, posizionato da Federica: challenge abbastanza
        alta (terza tacca, tra la terza e la quarta); più \emph{gameplay
        first}, ma con una componente narrativa importante, quindi un po'
        spostato verso la narrativa. Tra level design e combat, probabilmente
        un po' più verso il level design.
  \item \emph{Hollow Knight: Silksong} sarebbe più difficile e più
        \emph{gameplay first}: aumentano entrambe le dimensioni, ma il
        gameplay di più, forse anche per la quantità di contenuto rispetto a
        \emph{Hollow Knight}, e la narrativa è meno comprensibile. Rispetto
        al primo sarebbe spostato verso il combat.
  \item \emph{Castlevania: Aria of Sorrow}, secondo il docente, avrebbe una
        challenge forse un po' più bassa di \emph{Hollow Knight}.
  \item \emph{Metroid Dread}, più orientato allo \emph{stealth}, resta 2D
        come la gran parte dei titoli.\footnote{Nello stesso punto la
        trascrizione cita, come alternative 2D, titoli non riconoscibili
        (``Sintra, L'Enzion Mission, Tournus''); non si azzarda un'ipotesi.}
  \item La saga di \emph{Metroid Prime} va in prima persona; tutti gli altri
        in terza.
  \item Un titolo ``sistemico'' citato la volta precedente, con mira in
        prima persona e movimento in terza, andrebbe più verso la terza
        persona.\footnote{Il nome è illeggibile nella trascrizione
        (``Varen''). La descrizione (mira in prima persona, movimento in
        terza) corrisponde a \emph{Metroid: Other M}, che però compare poco
        dopo con un nome a sé (``Oderamma''); il riferimento resta quindi
        incerto.} Questi titoli avrebbero probabilmente una challenge più
        bassa degli altri: il primo \emph{Metroid} sta in un punto, e
        \emph{Metroid: Other M} in un altro, forse più in là.
  \item Gli studenti aggiungono altri titoli: \emph{Iconoclasts} (``bello''),
        \emph{Gato Roboto}, \emph{Animal Well} e alcuni altri.\footnote{Altri
        nomi proposti dagli studenti non sono riconoscibili nella
        trascrizione (``Rainbow Dash'', ``Pencil'', ``Mansions''); compare
        anche \emph{Firewatch}, citato come esempio di gioco in prima persona
        e 3D più che come metroidvania.}
\end{itemize}
Il docente lascia la mappa alla lavagna: chi vuole può aggiungere altri
giochi dalla lista ``non inclusiva'' di metroidvania, e qualcuno è invitato a
fotografarla e a mandarla anche a Mattia, collegato da remoto, che così la
può vedere meglio.

\paragraph{Che cosa mostra la mappa.} Qualcuno in aula nota che, guardando
la mappa, si vede anche \textbf{il genere che si preferisce}: tutti i giochi
più \emph{combat focused}, dice, li ha abbandonati. Il docente risponde che sulla
base di ciò che si inserisce e dei valori che gli si danno si riaggiusta
tutto il resto: su quella categoria, un titolo molto spinto sul combattimento
starebbe all'estremo \emph{combat focused}, e un altro
all'opposto.\footnote{I due titoli sono storpiati (``la InSource'' e ``il
regolamento''). Il primo potrebbe essere \emph{Nine Sols}, metroidvania
fortemente centrato sul combattimento, ma non è verificabile.}
Si vede anche come \textbf{gli elementi del mercato cambiano il valore}: se si
inserisce \emph{Silksong}, \emph{Hollow Knight} va riposizionato rispetto a
lui.

\paragraph{Assi polarizzanti.} I giochi risultano molto polarizzati, anche
perché alcuni assi scelti sono di per sé polarizzanti, come 2D e 3D: un gioco
in 2.5D come \emph{Dread} resterebbe comunque 2D. Si potrebbero definire
meglio le categorie: se il livello è a scorrimento orizzontale ma con
elementi in 3D, forse lo si potrebbe mettere a metà, ma diventa più
complesso. E già si vede che le zone sono quasi tutte popolate: ne sono stati
messi pochi; bisognerebbe inserirne un \textbf{grande numero} e cercare dove
ci sono i buchi, o combinazioni mai pensate. ``Più ne facciamo, più ne
troviamo.''

\paragraph{Non solo gameplay.} Gli assi possono essere anche artistici o
narrativi: lato artistico si potrebbe mettere il \textbf{realismo} da una
parte e un'arte più \textbf{stilizzata} dall'altra. Al corso interessa più il
lato gameplay, ma la perceptual mapping è un ottimo modo sia per fare
\textbf{ricerca} sia, potenzialmente, per fare \textbf{brainstorming}: trovare
le categorie che si vorrebbero avere e il modo di mescolarle per fare
qualcosa che non è mai stato fatto, o che si pensa non sia mai stato fatto
(si veda la tecnica dell'\emph{opposite direction},
sezione~\ref{sec:tecniche-idee}).

% ---------------------------------------------------------------------
\section{Altre fonti di dati}
\label{sec:fonti-dati}

Un altro modo è andare a cercare dati effettivi su piattaforme
specializzate. Il docente ne mostra una, limitata a Steam ma non l'unica, che
permette di vedere che cosa viene giocato e che cosa è di tendenza. Va capito
che cosa determina un trend: quali sono le categorie, quali i fattori. Ci
sono i \textbf{tag di Steam}, che spesso non si riescono a vedere o si fa
fatica a vedere, e i \textbf{dati di vendita}, ancora più difficili da
ottenere.

Scorrendo la pagina, il docente commenta qualche titolo: un gioco che non ha
ancora giocato e che i suoi amici vogliono giocare, e la classifica ``best of
2026'', in cui c'è un titolo in versione 1.0 che ``non avrei mai detto che
sarebbe finito qui dentro''.\footnote{Il titolo è storpiato nella
trascrizione (``Power 1.0''); potrebbe trattarsi di \emph{Palworld}, ma non
è verificabile. Anche il titolo citato subito dopo non è ricostruibile.} I
giochi si possono dividere per generi, oppure per gratuità, ``con dei grossi
asterischi'': \emph{Life is Strange} compare tra i gratuiti, ma ``è una
trappola'', perché gratuito è solo il primo episodio. Ci sono anche
categorie ``che forse è meglio non vedere''.

% ---------------------------------------------------------------------
\section{L'analisi degli utenti}
\label{sec:analisi-utenti}

\subsection{Uno dei punti più difficili}

Si passa al punto quattro della mappa, l'\textbf{analisi degli utenti},
già introdotta la volta precedente (sezione~\ref{sec:utenti-mercato}). Per il
docente è uno dei punti più difficili, certamente più dell'analisi dei
competitor, che diventerà progressivamente più facile con l'aumentare degli
strumenti a disposizione. L'analisi degli utenti richiede tanto lavoro di
ricerca e di domanda: creare un \emph{survey} apposito, con le domande
giuste, o anche solo chiedere feedback, richiede lavoro \textbf{prima}, per
decidere che cosa chiedere, e \textbf{dopo}, per prendere tutta la marea di
cose che sono state dette, segnarle, spacchettarle e capirle.

\subsection{Che cosa interessa sapere}

\paragraph{I dati quantitativi.} Degli utenti interessano anzitutto i dati
\textbf{quantitativi}: l'età, il genere (o ciò in cui la persona si
identifica), il lavoro, se ce l'hanno o no. Il lavoro aiuta a capire
\textbf{quanto tempo hanno a disposizione} e quanta ``potenza di fuoco''
economica. Se si vuole fare un gioco con numerose microtransazioni, magari
si deve fare un gioco \emph{mobile} in cui la \emph{stamina} si paga; e se ci
si vuole rivolgere a un target che lavora, conviene prediligere
\textbf{partite di durata più bassa}, ma la cui frequenza magari si paga, come
in \emph{Candy Crush}. Poi lo \textbf{status} e qualsiasi altro elemento utile
a descrivere ad alto livello chi è questo individuo: ``queste sarebbero le
\emph{game qualities} di una persona''.

\paragraph{Gusti e influenze.} Ci sono poi questioni più larghe. Età, sesso e
lavoro si descrivono in una parola o due; qui diventa più complesso. Che cosa
piace e che cosa non piace, non necessariamente in relazione al proprio
gioco: il tipo di musica che ascolta, il tipo di film che guarda, il rapporto
con anime e manga, che spesso sono una grande influenza sui videogiochi. Se si
prende uno stile specifico e lo si trasporta nel mondo videoludico, può
funzionare oppure no. L'esempio è quello fatto il lunedì precedente: un gioco,
di cui il docente non ricorda già più il nome, che prende a piene mani lo
stile dello \textbf{Studio Trigger} e lo riporta in ambito
videoludico.\footnote{Il titolo non è ricostruibile dalla trascrizione
(``il gioco di Destiny Beat''), e la lezione del lunedì a cui il docente fa
riferimento non è tra quelle trascritte nei capitoli precedenti.} Chi ama
quello stile ci va a nozze; altri un po' meno.

\paragraph{Che cosa fanno, e come: la routine.} Insieme a questo, è
importante capire la \textbf{routine}: come spendono la giornata. Sempre ad
altissimo livello: si sta cercando di delineare il target audience ideale del
proprio gioco. Quanto tempo hanno a disposizione; se giocano con altre
persone o da soli; se giocano in famiglia (una madre che gioca con i figli,
per esempio); quanto tempo dedicano al lavoro e quanto spazio resta per
giocare.

\subsection{Il caso Riot e League of Legends}

Una partita di \emph{League of Legends} dura in media, secondo il docente,
circa un'ora.\footnote{La durata media di una partita classificata di
\emph{League of Legends} è in genere più vicina alla mezz'ora; si riporta la
stima del docente, che non cambia il senso dell'esempio.} Se generalmente se
ne fa più di una (``non io: ho smesso con quella roba, con i tossici''), ci si
aspetta che il target audience ideale di quel gioco abbia \textbf{ampio
margine di tempo} da dedicargli.

\paragraph{Stessa IP, audience diversi.} \emph{League of Legends} esiste anche
per telefono, ma è un altro gioco: non è un \emph{porting} uno a uno della
partita.\footnote{Il titolo mobile è \emph{League of Legends: Wild Rift}.}
Sono due audience diversi. È un grande vantaggio, ed è il lavoro che ha fatto
Riot stessa (oggi sotto Tencent): prendere lo \textbf{zoccolo duro} del target
audience di \emph{League of Legends}, che è molto variegato, e offrirgli
esperienze diverse. La versione mobile ha una durata minore, per le persone
fidelizzate che vogliono quell'esperienza da MOBA anche mentre aspettano
l'autobus.

\paragraph{Gli altri giochi di Riot.} Il docente passa in rassegna gli altri
titoli di Riot, restando nel mondo di \emph{LoL} (quindi non
\emph{Valorant}):
\begin{itemize}
  \item un \textbf{picchiaduro}, con cui Riot ha cercato, secondo il docente
        senza successo, di convertire quell'utenza al genere dei
        picchiaduro.\footnote{Il titolo non è nominato in aula; si tratta
        quasi certamente di \emph{2XKO}. Si riporta il giudizio del docente
        sul suo esito così come è stato espresso.} Il margine era molto
        grande, perché le abilità di un giocatore di \emph{League of Legends}
        sono molto lontane da quelle richieste normalmente in \emph{Street
        Fighter}, \emph{Tekken}, \emph{Skullgirls} e negli altri picchiaduro
        tecnici; in più \emph{LoL} si gioca in team, mentre i picchiaduro
        sono generalmente uno contro uno;
  \item l'\textbf{auto battler}, ancora più mobile, accessibile e
        \emph{casual} (\emph{Teamfight Tactics});
  \item il \textbf{gioco di carte} di \emph{LoL} (\emph{Legends of
        Runeterra}; ``adesso non so in che stato è'');
  \item esperienze più \textbf{narrative}, come \emph{Song of Nunu}, una
        \emph{puzzle adventure} (``un'avventura platform 3D, mettiamola
        così'') con i protagonisti di \emph{LoL}. Prende un altro target, ma
        si assicura che una parte dell'utenza di base di \emph{League of
        Legends} finisca lì dentro.
\end{itemize}
Riot può permettersi di spostare persone perché nel tempo quel mondo si è
gonfiato moltissimo. Il docente stesso, quando giocava a \emph{League of
Legends}, era più interessato alla \emph{lore} del mondo che al MOBA in sé:
gli interessava leggere le storie e sentire i dialoghi tra i personaggi, su
cui a un certo punto Riot ha investito molto.

\paragraph{I due centesimi su Arcane.} Qui il docente si concede una
parentesi su \emph{Arcane}. Premette che sarà ``blasfemia'', ma lo stile
artistico ``alla \emph{Dishonored}'' non gli piace, e forse non è l'unico al
mondo.\footnote{Nella trascrizione i due nomi si sovrappongono (``Arcane'' e
``Arkane''): \emph{Arcane} è la serie animata di Riot prodotta dallo studio
Fortiche; Arkane Studios è lo sviluppatore di \emph{Dishonored}, il cui stile
pittorico viene spesso accostato a quello della serie.} Gli è piaciuto, nella
sua testa, fare tutti i collegamenti tra la lore della serie e quella del
gioco: a volte sono uno a uno, a volte si allineano, a volte no (Warwick,
per esempio, ``non c'entra nulla''). Il suo problema è che nella seconda
stagione sia stato ``ficcato a forza'' qualsiasi personaggio del mondo di
\emph{LoL}, quando non era necessario; e anche il fatto di rendere esplicita
la relazione tra Jinx e Vi, che nel gioco (``forse adesso lo è'') non era mai
stata dichiarata: c'era qualche tipo di legame, ma non si era mai detto che
fossero sorelle, come per Caitlyn. La prima parte ci può stare; la seconda è
caricata di troppi personaggi, e a lui ha pesato, perché il valore artistico
poteva essere più elevato. Non è \emph{fan service}: l'obiettivo finale è la
\textbf{trasmedialità} del prodotto, renderlo il più collegato possibile ad
altri elementi, ed è proprio scrittura transmediale, che si insegna. Non è un
problema di per sé: gli dà fastidio come utente, ma è ``uno su tutte le
persone'' che gli hanno dato cinque stelle (o una). ``Questi sono i miei due
centesimi.''

\paragraph{Spostare il target.} Il senso del discorso sul target è questo: se
si ha già un target di base, come Riot, e lo si ha molto frammentato, si può
pensare di \textbf{raggruppare e spostare} alcuni di quegli individui verso
altri giochi, altri prodotti, ``altri lidi''. Il modo più conveniente per
farlo, come si vedrà parlando di grafici, è il \textbf{clustering di dati};
ma non lo fa il singolo. È il lavoro del \emph{data analyst}, e servono una
mole enorme di dati e un algoritmo per separarli in modo sensato. Se ne
riparlerà, sfiorando l'argomento, perché non è comune che lo facciano gli
studenti, e non è comune nemmeno che le aziende medio-piccole abbiano un
sistema del genere.

\subsection{Chi gioca, chi giocherà}

Tornando alle slide: serve capire \textbf{chi gioca} il gioco, ma non solo:
chi lo giocherà, chi \emph{potrebbe} giocarlo, e in che spettro. Ci sarà il
\textbf{target audience dominante}, ma potrebbero esserci altre figure, altre
demografiche che gli ruotano attorno, potenzialmente interessate e che magari
si ha interesse a \textbf{convertire} in futuro: con i \textbf{DLC}, come detto
l'ultima volta, che sono già un modo per inserire meccaniche che potrebbero
piacere ad altri individui, o con gli \textbf{spin-off}, come appena visto.

\subsection{Il feedback online e l'ombra digitale}

Quali sono i loro feedback? Lo si è già visto per i competitor: se il proprio
gioco non esiste ancora, si va a vedere come reagisce il pubblico ai giochi
simili. Col passare del tempo diventa sempre più semplice: le reazioni sotto
gli streamer, i commenti su YouTube, Reddit.

\paragraph{Un grosso asterisco.} Ciò che si vede online, l'\textbf{ombra
digitale} (\emph{digital shadow}), non compone il 100\% dell'utenza.
\textbf{Chi è vocale online generalmente è polarizzato}: o ama alla follia un
prodotto o lo odia. Il docente non commenta sotto i video di YouTube, ma se
dovesse farlo andrebbe sotto il trailer di un nuovo \emph{Metroid} a dire
``figata, proprio quello che volevo''; altri andrebbero sotto lo stesso video
a chiedere ``fatemi una narrativa comprensibile''.\footnote{Il passaggio è
storpiato (``sotto il video di Ravenu [\dots] Samus che fa la merenda
grande''); il senso, un commento entusiasta a un nuovo capitolo di
\emph{Metroid}, è chiaro, il dettaglio no.} Sono due poli:
\begin{itemize}
  \item da una parte le persone che compreranno sicuramente il gioco, molto in
        target, e quelle che lo comprano e dicono agli altri di giocarlo: i
        cosiddetti \textbf{evangelisti}. Come un ex collega del docente che
        continua a dirgli ``compra \emph{Valheim}, giochiamo a
        \emph{Valheim}'': un evangelista, che cerca di convertirlo ``con
        scarsi risultati'';
  \item dall'altra le persone che il gioco non lo toccano, ed è giusto che
        non lo facciano, perché \textbf{nessun gioco è per tutti}, nessun
        prodotto è per tutti.
\end{itemize}

\paragraph{Minecraft è per tutti?} Uno studente obietta: \emph{Minecraft} non
è per tutti? Ha una base enorme, ma no, risponde il docente: come
\emph{Fortnite}, non è per tutti. \emph{Minecraft} è spintissimo verso
l'esplorazione e la comunità. Lo studente cita le \emph{Bed Wars}, che
trasformano il gioco in un PvP; il docente concede che lì c'è tanto, non solo
\emph{live ops} ma molto materiale della comunità: in \emph{Minecraft} si
può, e si fa, moddare il gioco, mentre in \emph{Fortnite} l'impatto che un
giocatore può avere creando contenuto è molto più ridotto. Resta il fatto che
\textbf{nel fulcro} \emph{Minecraft} è mirato all'esplorazione e alla
socialità, mentre \emph{Fortnite} alla socialità e alla competitività (lo si
rivedrà con le motivazioni, sezione~\ref{sec:quantic-foundry}). E restano la
stragrande maggioranza delle meccaniche giocate e percepite dai giocatori, e
quelle con cui il gioco si presenta nella prima o nelle prime due ore: quelle
sono, e quelle rimangono.

\subsection{Interviste, focus group e survey}

Identificato il target audience, in determinati punti del processo di
produzione ma anche prima, lo si \textbf{intervista}. Una volta fatta la
ricerca, gli si chiede, come si diceva l'ultima volta: ``stiamo facendo un
soulslike con i pony; ti piacerebbe? Perché no?''
(sezione~\ref{sec:utenti-mercato}). Serve come \textbf{cartina tornasole}: le
risposte del target audience confermano o smentiscono ciò che si è pensato
prima di metterlo a terra. Si chiedono opinioni: perché una cosa piace,
perché non piace.

\paragraph{I focus group.} Se si hanno i soldi si fanno \textbf{focus group}:
molto contenuti e omogenei, su diversi tipi di target audience, oppure misti,
con diversi tipi di utenti o demografiche. Poi i dati vanno analizzati.

\paragraph{Survey pagati e survey gratuiti.} Trovare le persone costa: chi
riceve un survey di questo tipo di solito viene pagato, poco ma pagato. Ci
sono piattaforme apposite, come \emph{SurveyMonkey}.\footnote{Il docente
aggiunge ``di Amazon, se non erro''. \emph{SurveyMonkey} è una piattaforma
indipendente; il servizio di Amazon per reclutare partecipanti pagati a
compiti e questionari è \emph{Amazon Mechanical Turk}. È probabile che i due
siano stati sovrapposti.} Il vantaggio è doppio: si ha la certezza che il
survey arrivi a certe persone e, utilissimo per un designer, che queste
\textbf{lo finiscano}, perché sono pagate per farlo e hanno un incentivo
molto elevato.

Se invece si è ``poveri'' (il docente si paragona all'immagine di un meme,
forse Mr.~Krabs, che non riesce a ritrovare), si fa il survey e lo si
``spamma'' il più possibile in giro. Da una parte si ottengono grandi entrate
di dati, che poi comunque vanno smazzate nell'analisi; dall'altra il survey
finisce davvero ovunque, e le persone che lo riempiono sono molto diverse tra
loro.

\paragraph{Le domande di filtro.} È un altro punto problematico, che il
docente aggiungerà alla lezione sui feedback. Con un survey diffuso così, una
parte del contenuto viene per forza ``mangiata'' dalle domande che servono a
separare i rispondenti: l'età, che tipo di giochi piace giocare, come si
spende il tempo, qual è il gioco preferito, e così via. Se si sta facendo un
metroidvania e la risposta alla domanda sul gioco preferito è ``non so,
\emph{The Sims}'', forse la persona non è proprio in target (per semplicità,
si assume di no). Per capirlo servono tante domande; ma \textbf{più domande ci
sono, meno la persona sarà incentivata a concludere} il survey.

\paragraph{La regola del pollice.} Il survey \textbf{non deve durare più di
cinque minuti}, e bisogna ``truffare'' le persone dicendo: ``questo survey dura
meno di cinque minuti, tre minuti per concluderlo, ti prego, in ginocchio,
completalo''. Se lo si paga, la persona può anche completarlo con calma;
altrimenti non lo farà.

\paragraph{Risposte a caso e outlier.} Il docente stesso, quando compila un
survey, alcune risposte le dà a caso; altre ``come per dire''; a volte,
quasi alla fine, lo interrompe e lo riprende dopo. Le singole risposte
perdono significato e diventano \emph{outlier}. Le soluzioni sono due:
\begin{itemize}
  \item avere davvero una marea enorme di entrate, \textbf{molto superiore al
        centinaio} (sezione~\ref{sec:dati});
  \item saper fare i survey, perché ``c'è un'arte'': più le domande sono
        semplici, limitate nelle risposte, puntuali, meno probabilità c'è
        che si risponda a caso. Si legge la domanda, la si capisce subito, si
        leggono le risposte, A, B, C, D, e si va dritti.
\end{itemize}
La tendenza si può ridurre, ma non ignorare del tutto, a meno di stare
dietro alla persona a guardare che cosa compila. Si può fare, per esempio
alla Games Week o a un'altra fiera\footnote{Nella trascrizione il secondo nome
è storpiato (``Fortale'', ``Force Playable''); non è ricostruibile.}: si fanno
le domande del survey stando ``sulla spalla'' e ci si assicura che lo si
compili. Ma anche così si modifica il modo in cui la persona risponde, che
potrebbe sentirsi indispettita: ``a ognuno il suo''.

\paragraph{Approfondimento: i consigli di Schell sui survey.} Nel capitolo sul
playtesting di \emph{The Art of Game Design}, Schell elenca alcuni consigli
che si sposano con quelli del docente: usare immagini quando si chiede di
elementi o scene del gioco; usare strumenti online come SurveyMonkey o Google
Forms, gratuiti o economici; evitare le scale da 1 a 10 e preferire
\textbf{scale a cinque punti} con ogni punto etichettato (da ``pessimo'' a
``eccellente''), che danno risultati più coerenti; \textbf{non mettere troppe
domande}, perché verso la fine le persone si distraggono e i risultati perdono
valore; somministrare il survey subito dopo la sessione di gioco; avere
qualcuno a disposizione per i chiarimenti; annotare età e genere di ciascun
rispondente; e soprattutto \textbf{non prendere i dati del survey come
vangelo}, perché difficilmente un survey è davvero scientifico e chi risponde
tende a inventare quando non è sicuro. Per le domande troppo complesse per un
modulo, Schell consiglia l'intervista, possibilmente privata e con una
traccia di domande pronta.

\subsection{Archetipi e osservazione}

``Questo discorso di fare survey è tosto'', e l'analisi dell'utenza è uno dei
punti più difficili: non se ne scappa. L'unico modo è prendere più dati
possibile e dedicare tanto tempo a smistarli. L'altro è creare il più
possibile \textbf{gruppi} o \textbf{archetipi}: approssimazioni di una serie
di utenti che però ``salvano la vita'' nella produzione, perché come sempre si
hanno tempo, soldi ed energie limitati. Se ne vedranno alcuni: uno molto noto
a breve, il più noto prima o dopo dicembre, a seconda del tempo.

\paragraph{Guardateli giocare.} Come già detto: si guardano interagire i
giocatori e, se possibile, li si \textbf{registra}, perché sono dati che poi
non vengono persi (sezione~\ref{sec:client-pitch}).

\paragraph{Approfondimento: nessuno aveva mai fatto un vero survey.}
Nell'introduzione di \emph{21st Century Game Design} (il libro di Bateman e
Boon) si racconta che l'industria ha prodotto per trent'anni giochi fatti da giovani
uomini occidentali per altri giovani uomini occidentali, saturando quel
mercato e lasciando scoperte altre demografiche. E perfino quel mercato, si nota,
non era davvero compreso: non tutti i giovani uomini sono uguali, preferiscono
sfide diverse, hanno quantità diverse di denaro e di tempo. Ernest W.~Adams,
autore dell'introduzione, scrive che in otto anni in Electronic Arts non ha
mai visto una ricerca di mercato davvero approfondita e condotta come si
deve: la comprensione dei giocatori si basava su supposizioni e
intuizioni. È il
problema che l'analisi degli utenti cerca di risolvere, ed è il punto di
partenza del loro modello demografico (sezione~\ref{sec:bateman}).

\subsection{La mappa dell'empatia}

Un altro modo di fare analisi degli utenti riguarda anche il \textbf{tono di
voce}: gli studenti lo vedranno, probabilmente, anche con Daniele. L'obiettivo
finale è sempre avere questo legame: \textbf{mettersi nei panni dell'utente}.
Per organizzare i dati rispetto a un certo target audience o archetipo di
utente, si inserisce tutto ciò che si è raccolto in \textbf{quattro
componenti} (tabella~\ref{tab:empatia}).\footnote{Lo strumento è comunemente
noto come \emph{empathy map} (mappa dell'empatia), con i quadranti
\emph{says}, \emph{does}, \emph{thinks} e \emph{feels}; il nome non viene
pronunciato in aula, ma le quattro componenti coincidono.}
\begin{description}
  \item[Che cosa dicono] Le citazioni, le cose dette esplicitamente.
  \item[Che cosa fanno] Le azioni. Si può dire ``sono il più grande giocatore
        horror dell'universo'', ma se davanti ad \emph{Amnesia} si gioca con
        le mani davanti agli occhi e la luce accesa, quello che si fa dice
        una cosa diversa.
  \item[Che cosa pensano] Le motivazioni, gli obiettivi. Si trovano anche con
        la ricerca, ma poi li si vede confermati: il ceto sociale di una
        persona, come si porta, lo status, il lavoro possono cozzare con ciò
        che pensa, con i suoi ideali, e così via.
  \item[Che cosa provano] Le emozioni, \emph{felt}: la lista di emozioni già
        vista (sezione~\ref{sec:vocabolario}). Qui si va nello specifico, nel
        tecnico: quanto più ci si indirizza già verso un modello o un
        framework, tanto più diventa utile, perché poi questi dati si
        riportano altrove.
\end{description}

\begin{table}[htbp]
  \centering
  \caption{Le quattro componenti per organizzare i dati su un utente o un
  archetipo.}
  \label{tab:empatia}
  \begin{tabular}{@{}p{0.2\textwidth}p{0.32\textwidth}p{0.4\textwidth}@{}}
    \toprule
    \textbf{Componente} & \textbf{Che cosa raccoglie} & \textbf{Esempio dalla lezione} \\
    \midrule
    Dicono (\emph{says}) & citazioni, affermazioni esplicite & ``sono il più grande giocatore horror'' \\
    Fanno (\emph{does}) & azioni, comportamenti osservati & gioca ad \emph{Amnesia} con le mani sugli occhi \\
    Pensano (\emph{thinks}) & motivazioni, obiettivi, ideali & status e lavoro che cozzano con gli ideali \\
    Provano (\emph{feels}) & emozioni & la lista di emozioni del corso \\
    \bottomrule
  \end{tabular}
\end{table}

\paragraph{Senza empatia non c'è design.} L'analisi degli utenti è
incredibilmente importante perché dà i dati necessari per
\textbf{empatizzare} con l'audience. Non significa doverne vestire sempre i
panni; ma in certi momenti di discussione e di design è utile avere questa
``maschera'' da indossare (sezione~\ref{sec:utenti-mercato}).
\textbf{Senza empatia non c'è design}: l'empatia è uno dei punti di arrivo.

\paragraph{Approfondimento: la lente del giocatore.} Schell formalizza lo
stesso atteggiamento nella \emph{Lens of the Player} (lente n.~19): smettere
di pensare al proprio gioco e cominciare a pensare al proprio giocatore. Le
domande sono: in generale, che cosa gli piace? Che cosa non gli piace, e
perché? Che cosa si aspetta di vedere in un gioco? Se fossi al suo posto, che
cosa vorrei vedere? Che cosa gli piacerebbe, o non gli piacerebbe, del mio
gioco in particolare? Un buon designer, scrive Schell, dovrebbe essere sempre
un \emph{avvocato del giocatore}.

% ---------------------------------------------------------------------
\section{L'analisi di mercato}
\label{sec:analisi-mercato}

\subsection{C'è spazio per il mio gioco?}

L'analisi di mercato estende ciò che si è visto con la perceptual mapping e
con il discorso sull'audience (sezione~\ref{sec:utenti-mercato}). ``Mercato''
può voler dire un continente, un paese, un mercato saturo o non saturo. La
domanda è: \textbf{esiste spazio per il mio gioco in questo momento storico?}
Quanto è grande il mercato? È una \textbf{nicchia}? Come detto l'ultima volta,
i giochi alla \emph{Papers, Please} sono una nicchia; un altro genere, ormai,
è un mercato molto ampio ma anche molto saturo.\footnote{Il genere è storpiato
nella trascrizione (``Solo Stripe''); dal seguito della lezione potrebbe
trattarsi dei soulslike, ma non è certo.}

\paragraph{I giochi brevi.} Uno studente fa notare che il mercato dei giochi
simili a \emph{Papers, Please} non è saturo: ogni volta che ne esce uno, sono
giochi piccoli, la gente se li gioca, e poi non ce n'è più. Il docente
apprezza: è un discorso molto interessante. Molti di questi giochi, come
anche molti giochi mobile, hanno una vita o una durata molto ridotta, e
permettono all'utente di provarne e finirne molti in poco tempo; cosa che si
fa meno con i AAA e i AA, soprattutto con gli open world.

\paragraph{Monopolizzare l'utente.} Quando si presenta un gioco AAA,
la \textbf{longevità} del prodotto deve essere un punto di forza. Ci sono
giochi che puntano il più possibile a \textbf{monopolizzare l'utente}: se un
open world ha cento livelli, mille espansioni, ``duecento milioni di biomi'',
il giocatore che gioca a quello non gioca all'open world B. Vale anche per gli
\emph{hero shooter}: chi è fidelizzato su \emph{Valorant} o su
\emph{Counter-Strike} tendenzialmente non andrà a giocare a \emph{Overwatch}.
È anche uno dei motivi, oltre a una serie di altri, per cui \emph{Concord} ha
fallito.

\subsection{Seguire i trend}

È discutibile se seguire i trend sia un modo corretto di procedere, anche
perché \textbf{i trend si seguono sempre in ritardo}: anticiparli è
difficilissimo, e fanno fatica anche le aziende AAA.

\paragraph{I casi del pubblico dominio.} Un caso in cui il trend era
prevedibile è quello delle opere che diventano libere da diritti. È successo,
in ambito cinematografico, con lo \emph{Steamboat Willie} originale (non
Topolino in generale, ma quella prima versione) e con \emph{Winnie the Pooh}:
quando sono scaduti i diritti d'autore, chiunque ha potuto usarli in qualsiasi
prodotto.\footnote{Il docente dice che queste opere sono ``entrate in Creative
Commons''; più precisamente sono entrate nel \textbf{pubblico dominio}, per
scadenza del diritto d'autore (negli Stati Uniti nel 2022 per il primo libro
di \emph{Winnie the Pooh} e nel 2024 per \emph{Steamboat Willie}). Le licenze
Creative Commons sono invece concesse volontariamente dagli autori.} È
successo anche con Lovecraft e Cthulhu: c'è stata una valanga di giochi,
titoli, scritti e film lovecraftiani. Molti prodotti hanno fatto la
produzione prima e poi, alla scadenza, hanno rilasciato subito, tutti insieme
o comunque nel tempo: c'è stata una bella finestra temporale in cui molti
hanno usato temi cthulhiani.

\paragraph{È difficile prevedere il mercato.} Probabilmente Diego, il
direttore del corso, direbbe che una certa grande azienda è brava a farlo; è
vero, ma anche lei di recente ha fatto ``dei grossi
inciampi''.\footnote{Il nome dell'azienda non è ricostruibile dalla
trascrizione (``Without Bordecost'', ``sotto Asgore'').} Il punto è che fare
tanta ricerca, e farla bene, può assicurare una buona fetta di mercato e in
generale la \textbf{sopravvivenza dell'azienda}; farla male può causare il
contrario, il fallimento. E spesso il fallimento arriva anche con ricerche
mediocri, o perfino positive, per una serie di altre questioni, prevedibili o
meno, come gli andamenti del mercato, che sono poco prevedibili. Bisogna
dedicarci tempo e attenzione.

\subsection{Lo stato emotivo del mercato}

Qual è lo stato ``emotivo'' del mercato, ``mettiamola così''? Gli utenti sono
contenti di ciò che c'è?
\begin{itemize}
  \item Il mercato dei \textbf{soulslike} non solo è saturo, ma è saturo di
        titoli che la gente gioca con \textbf{moderato apprezzamento}. C'è
        sempre la ``grande torta cicciotta'' di FromSoftware, che non ne
        sbaglia quasi una, e poi una serie di ``cupcake piccolini'', tutti AA,
        ``tripla I'' o indie, con questo tipo di meccaniche.
  \item Il mercato dei \textbf{monster collectible} è il contrario: non è
        saturo e, in più, gli utenti sono scontenti.
\end{itemize}

\subsection{Barriere, tasse e logistica}

Ci sono barriere? Sono quelle viste l'ultima volta
(sezione~\ref{sec:utenti-mercato}): si può o non si può fare qualcosa in quel
mercato? Conviene pensarci prima, ed eventualmente decidere
\textbf{consapevolmente} se affacciarsi o no.
\begin{itemize}
  \item Se il gioco è completamente basato sulle \textbf{loot box}, si
        potrebbe non poterlo lanciare in \textbf{Belgio}: \emph{Overwatch} ha
        tolto quel sistema in quel paese, e molti giochi mobile in Belgio non
        escono proprio per questa ragione.
  \item Ci sono \textbf{tasse}, regolamenti, regole. \emph{Kickstarter}, per
        esempio, trattiene una percentuale sul titolo che si sta sviluppando,
        e quella percentuale va considerata.
  \item Se si fa un \textbf{gioco fisico}, bisogna considerare non solo i
        costi di produzione, ma il numero di copie, quante ne vanno vendute
        per avere un ritorno, e dove vanno stoccate: ``non è che le mettete
        nel giardino di casa'', serve un \textbf{magazzino}. E poi chi si
        occupa dei trasporti, chi dell'imballaggio, e così via.
\end{itemize}

Il docente valuta se fare la pausa subito o dopo, perché la parte successiva
è ``molto importante'', e decide per subito: cinque minuti di
pausa.\footnote{Durante la pausa la trascrizione automatica registra solo
frasi ripetute in loop, prive di contenuto, che qui si omettono.}

% ---------------------------------------------------------------------
\section{L'analisi di gameplay: le motivazioni di Quantic Foundry}
\label{sec:quantic-foundry}

\subsection{Un modo, non l'unico}

Dopo la pausa si passa all'\textbf{analisi di gameplay}. Il docente precisa
che quella che presenta non è l'unica: è uno dei modi possibili, ed è
importante, perché ritornerà in futuro. Segnala anche una serie di talk che
si possono cominciare a vedere, anche se non sa se consigliarli caldamente.

\paragraph{Il Dragon Speech.} Prima, un'altra raccomandazione che gli è tornata
in mente: il \textbf{Dragon Speech} di Chris Crawford. ``Dovete andare a
vederlo, non lo ripeto più. È importante.''\footnote{Si tratta del celebre
intervento di Chris Crawford alla Computer Game Developers Conference del
1992, in cui il designer annunciava il proprio abbandono dell'industria
paragonando la propria ricerca a quella di un drago.}

\subsection{Quantic Foundry}

Il docente apre il sito di \textbf{Quantic Foundry} e nota che è stato
rifatto: ora è ``vestito'' in modo più carino, accessibile, giocoso e saturo di
colori; a lui piaceva di più quello vecchio, in bianco e nero, ``preciso,
elegante'', più stoico, ``grezzo e duro, spartano''.

Quantic Foundry ha raccolto dati da un numero enorme di giocatori (il
docente parla di circa due milioni), e i suoi fondatori, uno dei quali è un
\emph{data analyst}, hanno cercato di correlarli per arrivare a una serie di
risultati interessanti.\footnote{I fondatori di Quantic Foundry sono Nick Yee
e Nicolas Ducheneaut; nella trascrizione il primo compare come ``Nick Lie''.}
Alcuni li si vedrà più a fondo, come gli \textbf{archetipi di giocatori} di
Quantic Foundry, che sono una cosa solo loro. Ciò che interessa ora è uno
\textbf{spettro}: categorizzare le necessità e le motivazioni dell'utente
secondo le emozioni della lista già vista (sezione~\ref{sec:vocabolario}).

\paragraph{I limiti dei dati.} I dati sono \textbf{chiusi}: non possono essere
visti se non da soci associati. E sono limitati ai \textbf{giocatori
maggiorenni}: c'è sempre il limite del non fare ricerca statistica su
minorenni, o comunque su giocatori molto piccoli. La ``personalissima teoria
del complotto'' del docente (da raccontare per far ridere e dimenticare
subito dopo) è che Yee abbia sempre avuto il pallino di estendere la ricerca
anche ai più piccoli; ma servirebbe il consenso dei genitori, e le
regolamentazioni cambiano da stato a stato: ``è un casino''. Il docente non
ha idea, per esempio, di come abbiano gestito i dati raccolti in Russia
rispetto a quelli in America, dove la soglia della maggiore età è diversa.

\paragraph{Il talk.} Il GDC Talk su queste \textbf{dodici motivazioni} si trova
sul sito della GDC e anche su YouTube (cercando ``Quantic Foundry talk
GDC''). Ce ne sono altri degli stessi autori; il docente consiglia solo il
primo, se proprio lo si vuole vedere, ma se ne riparlerà comunque.

\subsection{Sei gruppi, dodici motivazioni}

Le motivazioni sono raggruppate a due a due in sei categorie, che per ora il
docente chiama ``emozioni''. Ciascuna ha un corrispettivo nella lista di
emozioni del corso (sezione~\ref{sec:vocabolario}):
\begin{description}
  \item[Action] corrisponde all'\emph{excitement}.
  \item[Social] corrisponde a \emph{socialization}, \emph{fellowship},
        \emph{community}, ``chiamatela come volete''.
  \item[Mastery] è la padronanza, la \emph{challenge}; ma è proprio lo
        sviluppo delle skill legate all'azione del giocatore, quello visto
        nel framework di Bura (sezione~\ref{sec:bura}).
  \item[Achievement] corrisponde all'\emph{accomplishment}: ``è la stessa
        roba''.
  \item[Immersion] corrisponde alla \emph{fantasy}, ed è anche
        \emph{vicarious experience} quando l'esperienza è molto lontana dalla
        nostra: non il pilota di corse, ma il cavalcatore di draghi.
  \item[Creativity] corrisponde alla \emph{creation}: poter giocare con le
        proprie idee, dare sfogo alla propria immaginazione tramite gli
        strumenti del gioco, come creare il Taj Mahal in \emph{Minecraft}.
\end{description}
Ciascuna categoria contiene \textbf{due motivazioni}, e ciascuna motivazione
è \textbf{uno spettro}: non è che un gioco ha o non ha la community; si
colloca su uno spettro (tabella~\ref{tab:quantic}).

\begin{table}[htbp]
  \centering
  \caption{Il modello delle motivazioni di Quantic Foundry, con i poli e gli
  esempi discussi in aula.}
  \label{tab:quantic}
  \small
  \begin{tabular}{@{}p{0.13\textwidth}p{0.15\textwidth}p{0.32\textwidth}p{0.32\textwidth}@{}}
    \toprule
    \textbf{Gruppo} & \textbf{Motivazione} & \textbf{Polo alto} & \textbf{Polo basso} \\
    \midrule
    Social & Community & cooperazione: \emph{WoW}, \emph{LoL}, \emph{Destiny}, \emph{Rainbow Six} & indipendenza: giochi LEGO, \emph{Hatoful Boyfriend} \\
    & Competition & alto conflitto: \emph{LoL}, \emph{Dota 2}, \emph{CoD}, \emph{CS:GO}, \emph{For Honor} & basso conflitto: \emph{Firewatch}, \emph{Myst}, \emph{Gone Home} \\
    Action & Excitement & \emph{thrilling}: \emph{CS:GO}, \emph{Ready or Not}, \emph{Bodycam}, \emph{Smash Bros.} & rilassato: \emph{Pokémon}, \emph{Civilization}, \emph{Animal Crossing} \\
    & Destruction & caos: esplosioni, armi, \emph{Doom}, muri che crollano & calma: \emph{Professor Layton}, \emph{Harvest Moon} \\
    Achievement & Completion & \emph{task oriented}: JRPG, \emph{The Witcher}, \emph{Zelda} & \emph{self-driven}: \emph{Minecraft}, \emph{Kerbal Space Program} \\
    & Power & progressione di potere: \emph{Dark Souls}, \emph{Diablo 3}, \emph{Borderlands} & efficienza: \emph{Mario Odyssey}, \emph{Night in the Woods}, \emph{Ico} \\
    Mastery & Strategy & contemplativo: \emph{Europa Universalis}, \emph{Balatro} & spontaneo: \emph{Candy Crush}, match-3, \emph{Mario Kart} \\
    & Challenge & \emph{skill based}: \emph{Dark Souls}, \emph{Hollow Knight}, \emph{Ori} & \emph{easy fun}: \emph{Oxenfree}, \emph{Animal Crossing} \\
    Immersion & Fantasy & lore profonda: \emph{Mass Effect}, \emph{Dark Souls}, \emph{Fallout} & setting generico: \emph{Street Fighter}, \emph{Candy Crush} \\
    & Story & dramma: Telltale, BioWare, \emph{Persona} & \emph{open ended}: \emph{Factorio}, \emph{SimCity} \\
    Creativity & Discovery & esplorazione: \emph{Skyrim}, \emph{Minecraft}, \emph{Subnautica} & pratico: \emph{FIFA}, \emph{Madden}, scacchi \\
    & Design & \emph{self-expression}: personalizzazione del personaggio & personaggio fisso: \emph{Braid}, \emph{Spelunky} \\
    \bottomrule
  \end{tabular}
\end{table}

\subsection{Community}

Il primo spettro è la \textbf{community}: condividere esperienze. Da una
parte c'è tutto ciò che è \textbf{collaborativo}: non solo incontrare altri
esseri umani, ma interagirci insieme. Dall'altra il gioco completamente
\textbf{single player}, in cui il giocatore è in controllo della propria
esperienza: cooperazione da una parte, \textbf{indipendenza} dall'altra.

\paragraph{Una parentesi: la community fuori dal gioco.} Un gioco
completamente single player può avere zero community \emph{dentro} il gioco.
Ma esiste la community di \emph{Hollow Knight} fuori dal gioco: prima si
parlava del \emph{meta} delle build di \emph{Hollow Knight}, una cosa
percepibile solo affacciandosi fuori. Se ne riparlerà.

\paragraph{Gli esempi.} Giochi con il \emph{teamwork} alle stelle: qualsiasi
\emph{Warcraft}, potenzialmente anche \emph{League of Legends},
\emph{Destiny}, \emph{Rainbow Six}. Giochi single player: i giochi LEGO, di
solito, e \emph{Hatoful Boyfriend}, il \emph{dating sim} con i piccioni
(``incredibile che ci sia''). A metà ci sta \emph{Stardew Valley}: non è
un'esperienza solo single player, ha la possibilità di cooperativa, ma la coop
non è così spiccata e i suoi sistemi non sono così spinti in quella
direzione.

\paragraph{Stardew Valley contro Animal Crossing.} Uno studente propone
\emph{Animal Crossing} come gioco con un discorso di community. Il docente è
d'accordo che sarebbe più spinto da quella parte rispetto a \emph{Stardew
Valley}. In \emph{Stardew Valley} si condivide la propria fattoria con le
persone che giocano con te: per ogni slot di salvataggio ci sono quattro
posti; uno è il tuo, e una volta creati i personaggi per gli altri tre, per
quanto ne sa il docente, non li si può più cancellare. In \emph{Animal
Crossing} si può andare a visitare l'isola di un altro e farlo venire sulla
propria; le meccaniche sono semplici, e comunque qui si parla di motivazione,
quindi di responso emotivo. Il docente è dell'idea che \emph{Animal Crossing}
sia più verso il teamwork, ma ``andrebbero trovati i dati''.

\subsection{Competition}

La \textbf{competition} è competere con un altro individuo, generalmente con
una risoluzione uno a zero: uno vince, l'altro perde. Da una parte ci sono i
giochi ad \textbf{alto conflitto}: competizioni dirette, arene, duelli, match
uno contro l'altro, eventualmente \emph{leaderboard} che mettono in rapporto
ciò che si è fatto con ciò che hanno fatto gli altri, e \emph{ranking}, cioè
una scala gerarchica.
\begin{itemize}
  \item \emph{League of Legends} è ad alto conflitto: ha i ranghi, ``da
        bronzo, legno, mattone fino a diamante incredibile, mega ultra top
        1000 mondiale''. \emph{Dota 2} è lo stesso; \emph{Call of Duty} e
        \emph{CS:GO} hanno la stessa cosa.
  \item Anche i \textbf{giochi di carte} sono generalmente ad alto conflitto.
\end{itemize}

\paragraph{Reattivo o strategico, non importa.} L'alto conflitto
\textbf{non è legato al tipo di gameplay}, reattivo o strategico (se ne
parlerà). In \emph{Call of Duty} c'è la reazione \emph{second to second}:
scegliere dove andare, mirare, reagire. In \emph{Hearthstone} si ha il proprio
turno, circa un minuto, da giocare in tranquillità, anche fino agli ultimi
dieci secondi: si ha modo di pensare alle proprie azioni. Ma il risultato,
lato competizione, è lo stesso.

\paragraph{Dark Souls e For Honor.} \emph{Dark Souls} non è di per sé un gioco
ad alto conflitto; ma se ci si concentrasse solo sul multiplayer, quella
sfera lo sarebbe. Un esempio di alto conflitto ``puro'' è \emph{For Honor} di
Ubisoft, l'arena di combattimento medievale con classi come il vichingo e il
cavaliere (il docente lo cerca su YouTube, tra pubblicità di ogni genere, e
ne mostra un combattimento). Ci sono anche dei \emph{minion}, come in
\emph{League of Legends}, NPC con poca vita, ``carne da macello''; ma nei
duelli uno contro uno c'è il \emph{target lock} e un sistema di guardia molto
interessante basato sulle \textbf{stance}: la spada in basso a destra, in
basso a sinistra, in alto, e si difende dalla parte da cui arriva il colpo.
Funziona come un ``carta, forbice, sasso'': se si entra dove l'avversario è
scoperto si fanno più danni. L'obiettivo è parare quando arrivano i colpi e
farli entrare dove l'altro non è difeso. Un sistema che, secondo il docente,
si potrebbe ``rubare'' e mettere altrove per fare un gioco interessante.

\paragraph{Il basso conflitto.} Dall'altra parte non c'è un avversario: niente
duelli, niente sfida contro un altro giocatore. \emph{Firewatch},
\emph{Myst}, \emph{Gone Home}, i \emph{walking simulator} in generale: non si
è in competizione né con sé stessi né con altri. Anche i giochi LEGO, a
bassa competizione per definizione, e \emph{Dragon Age 2}.

\paragraph{E le mod?} Uno studente chiede che cosa succede se una mod non
ufficiale aggiunge del competitivo a un gioco che non lo ha.\footnote{Nella
trascrizione il docente si rivolge a ``Valeria''; nel resto della lezione
compare uno studente di nome Valerio, ma non è possibile stabilire con
certezza chi abbia fatto la domanda.} Bisognerebbe capire \textbf{quanto
impatta}: magari la mod risponde a un'esigenza dell'utenza (``vorrei il gioco
più competitivo''), ma vale lo stesso discorso di \emph{Dark Souls}: la
stragrande maggioranza del gioco è single player, è pensata come tale, e la
sfera della competizione è molto ridotta. Il gioco si sposta sulla linea, ma
non poi così tanto. Su ogni spettro \textbf{ci si sposta aggiungendo o
togliendo contenuto} in una direzione specifica; va capito di quanto.

\paragraph{Le motivazioni dominanti.} Ogni gioco ha \textbf{ciascuno} di
questi elementi, più o meno spinto. Per semplicità, come si fa con le
emozioni, per descrivere un gioco si prendono gli elementi
\textbf{dominanti}: le motivazioni più spinte, le più identificative. Per
\emph{League of Legends}, tra \emph{action}, \emph{social}, \emph{mastery},
\emph{achievement}, \emph{immersion} e \emph{creativity}, probabilmente
interessano \emph{mastery}, \emph{social} e \emph{action}, in quest'ordine.
Significa che non ha \emph{achievement}, o \emph{creativity}? No: li ha, ma in
misura molto ridotta rispetto al resto. ``Questo è il concetto.''

\subsection{Excitement}

L'\textbf{excitement} è quanto il gioco è \emph{fast paced}, quanto è
\emph{action}, quanti picchi di adrenalina ci sono in determinati secondi di
gioco, nel \emph{second to second}. Gli FPS in prima persona come \emph{CS:GO}
sono \emph{thrilling}; ancora più thrilling potrebbe essere \emph{Ready or
Not}, e ancora più di questo \emph{Bodycam}, che spinge il realismo più in
alto. Ma vale anche per \emph{Smash Bros.}, chiaramente \emph{fast paced},
\emph{action based}, \emph{adrenaline rush}; e anche i giochi di corse sono
generalmente thrilling, compreso il nuovo gioco di corse di \emph{Star Wars}.

Dall'altra parte ci sono i giochi rilassati: qui si vedono già i due grandi
tipi di gameplay che si studieranno, \textbf{action reactive} contro
\textbf{strategico}, \emph{planning} e \emph{predictive}; o, come si dice di
solito, \emph{pace} elevato contro \emph{pace} lento e calmo. Ci sono tutti i
\textbf{turn based}, che possono essere messi in pausa senza conseguenze e
tendono al rilassamento e alla prevedibilità: \emph{Pokémon} ci ricade.
Bassa stimolazione visiva (non all'estremo) come \emph{Myst};
\emph{Civilization}, che non è all'estremo perché a seconda dell'azione un po'
di thrilling c'è, ma il fatto che sia a turni lo rende molto più calmo.
\emph{Animal Crossing}, chiaramente, è tutto spinto in questa direzione.

\paragraph{Smash è excitement o competition?} Uno studente chiede se
\emph{Smash Bros.} non vada piuttosto sotto la competizione. Il docente è
d'accordo: sicuramente c'è il thrilling, ma se si dovesse scegliere la
motivazione principale, probabilmente avrebbe ragione lo studente. Sono i
giochi \textbf{horror} a essere ``questa roba'': di un \emph{SOMA} o di un
\emph{Amnesia} è più necessario dire che sono thrilling. Per gli sport, la
competizione è la motivazione principale.

Lo studente gioca a \emph{Smash} in modo competitivo, e fa quindi parte del
target audience; ma ci sono persone che giocano a \emph{Smash Bros.} come
\emph{party game}, ed è probabilmente il suo punto di forza. Lo studente
ribatte che \emph{Ultimate} è stato pensato anche come party game, ma che
alcune aggiunte di meccaniche lo hanno reso più difficile; il docente
concede che c'è anche quella parte.

\paragraph{Mario Kart.} Un altro esempio Nintendo: in \emph{Mario Kart} la
competizione è un punto forte, ma meno rilevante dell'excitement, perché è
più un party game. L'obiettivo non è tanto vincere quanto partecipare: è più
interessante vincere tutte le corse, oppure, anche arrivando ultimi,
assicurarsi che l'amico con cui si gioca si becchi il guscio blu prima della
linea del traguardo? Si dice che \emph{Mario Kart} distrugga le amicizie, ed è
pensato per quello. Il fatto che il risultato si possa ribaltare anche a un
centimetro dalla fine è pensato proprio per rendere la competizione molto più
leggera e il gioco molto più \emph{party} (si ragionerà anche sul suo
\emph{loop} di gioco; e dipende dal \emph{Mario Kart}).

\subsection{Destruction}

La \textbf{destruction} è l'entropia, la possibilità di modificare e
distruggere l'ambiente. Ricade sempre nell'\emph{action}, perché è la
necessità del giocatore di ``fare casino'': distruggere cose, impattare con
violenza il mondo di gioco o le persone che ci stanno dentro. Da una parte,
nel reame del \textbf{caos}, ci sono le esplosioni, le pistole, il sangue, la
violenza: strappare l'occhio al capo demone, far scoppiare qualcosa in un
gioco di guerriglia e vedere il muro crollare fisicamente.

Dall'altra ci sono i giochi il cui mondo non si può distruggere. In
\emph{Professor Layton} non si può distruggere letteralmente niente: il mondo
è fisso, si gira per la città (o per il treno, a seconda del capitolo), al
massimo si parla con le persone. Non ci sono armi, non c'è \emph{gore}, e
generalmente nessuno muore (in qualche capitolo forse si vede ``una mano'', e
basta). Ci ricadono anche \emph{Animal Crossing} e \emph{Harvest Moon};
\emph{Stardew Valley} sta un gradino sopra, perché si possono fare piccoli
danni: c'è la bomba, si possono distruggere i \emph{tile} piazzati da sé, i
mostri, anche i simpatici slime con il sorriso. Ma una bomba piazzata accanto
a un NPC non gli fa nulla, né distrugge la sua casa.

\subsection{Completion}

La \textbf{completion} è avere o non avere \textbf{linee guida dirette}, vere
e proprie \emph{checklist} da riempire. Molti JRPG, molti \emph{Dragon Age},
in un certo senso anche \emph{The Witcher}, hanno obiettivi primari e
secondari, \emph{quest} principali e secondarie; anche \emph{Zelda} è così.
Diventa ancora più spinto tutto ciò che è legato agli \emph{achievement}:
``batti il boss dieci volte per sbloccare questa cosa'' è \textbf{task
oriented}. Dall'altra parte c'è il \textbf{self-driven}: è il giocatore a
decidere le proprie priorità. \emph{Minecraft} sta qui, come generalmente
tutti i \emph{sandbox}. \emph{Kerbal Space Program} è questo: bisogna mandare
gli astronauti alieni nello spazio, ma si decide come montare il razzo, nei
limiti del gioco.

\subsection{Power}

Il \textbf{power} è un \textbf{falso amico}, e lo si capirà parlando di
progressione. Indica la possibilità di \textbf{ottenere potere}: livellare,
cioè modificare le statistiche, e migliorare le armi. Qui c'è sicuramente
\emph{Dark Souls}; e tutti i giochi come \emph{Diablo 3} o \emph{Borderlands},
con N difficoltà che scalano, in cui il giocatore ``lutta'' per ottenere la
stessa arma, con la stessa forma e probabilmente la stessa rarità, ma con un
numero N volte più alto. Tutto ciò che si può migliorare tramite statistiche
o tramite lo sblocco di nuove meccaniche.

Il contrario è la \textbf{self-efficiency}, che il docente aveva anticipato
nell'esercizio sui metroidvania: in \emph{Super Mario Odyssey} non si diventa
più forti. Ci sono nemici che, catturati, sbloccano nuovi verbi di gioco, ma
sono limitati a dove si trova il nemico. Nella parabola di \emph{Mario
Odyssey} si sviluppa \textbf{competenza nel platforming}: si capisce meglio
come lanciare il cappello, come fare il doppio salto sul cappello, e così via.
Lo stesso in \emph{Super Mario 64}.

All'estremo assoluto ci sarebbe un personaggio che ha già tutti i verbi di
gioco e tutte le statistiche al massimo, o meglio, le cui statistiche
\textbf{non cambieranno mai}. \emph{Night in the Woods}, avventura narrativa, è
così. \emph{Firewatch} no, perché si sbloccano cose, quindi sta un po' più in
là, nonostante sia spiccatamente narrativo. \emph{Her Story} è così: nessuna
progressione sul personaggio.\footnote{Nella trascrizione il titolo compare
come ``Air Story''; la lettura \emph{Her Story} è la più plausibile. Subito
dopo il docente cita un altro titolo che dice di non aver giocato (``Easy
Days''), non riconoscibile.} \emph{Ico} è così: c'è il protagonista e la
ragazza da portare via tenendola per mano; ``quello è, quello rimane'', e il
protagonista non sblocca abilità.

\subsection{Strategy}

La \textbf{strategy} è la complessità nelle decisioni: quanto il gameplay è
strategico, o il contrario. Lo si vedrà definito in altri modi, ma in
generale un gioco con molta excitement avrà poca strategia in questo senso;
non è sempre detto, dipende dal framework. Da una parte c'è il
\textbf{contemplativo}: poter pianificare, avere decisioni complesse e
strategie a lungo termine, e il \textbf{tempo} per farle. \emph{StarCraft}
non sta su questa linea, perché c'è comunque un gameplay \emph{second to
second} a cui sottostare: il docente mostra un video dei migliori giocatori,
che cliccano velocissimi (``ho cliccato troppe volte, scusate'').
\emph{Europa Universalis} è ancora più contemplativo; \emph{Balatro} è
completamente strategico: ci sono i turni, nessun limite di tempo, e bisogna
scegliere che scala fare, se rischiare o no, se cambiare mano o no.

Dall'altra parte c'è lo \textbf{spontaneo}: gameplay reattivo e \textbf{basso
peso cognitivo}. È diverso dall'excitement, che spesso ha un alto peso
cognitivo: qui si fanno azioni in poco tempo con un basso costo cognitivo.
\emph{Disney Emoji Blitz}, \emph{Candy Crush}, i \emph{match-3}; e \emph{Mario
Kart}, con il suo gameplay reattivo nel \emph{second to second}: si deve
capire quando fare la curva, ma la si fa subito.

\subsection{Challenge}

La \textbf{challenge}, la sfida, va da \textbf{skill based} a \textbf{easy
fun}. Il docente precisa subito che non è un errore: è probabilmente
l'\textbf{unica eccezione} per cui si può usare il termine \emph{fun}, bandito
dal vocabolario del corso (sezione~\ref{sec:vocabolario}). Sono i
\textbf{quattro tipi di fun di Nicole Lazzaro}: \emph{easy fun}, \emph{hard
fun}, \emph{serious fun} e \emph{people fun}, ``i quattro tipi di fun che
potete citare''. Il \emph{people fun} è, per esempio, la community; l'\emph{hard
fun} è letteralmente il contrario dell'\emph{easy fun}, cioè la sfida. Si
rivedranno più avanti.

\paragraph{Skill based.} \emph{Dark Souls}, ma ancora di più, come si diceva
prima, \emph{Ori and the Blind Forest}: è skill based e anche molto
reattivo. Si pensi alle sequenze di fuga, o al fatto che si salta
``rimbalzando'' sui proiettili dei nemici; c'è del \emph{freestyle}. Ci sono
\emph{forgiveness mechanics} molto importanti (``ci mancherebbe''), ma è skill
based. Lo è anche \emph{Hollow Knight}: messi in ordine, ci sarebbe prima
\emph{Hollow Knight} e poi \emph{Ori}, anche se secondo il docente non sono
così lontani. Il docente ha giocato solo il primo \emph{Ori}.

\paragraph{Una digressione su Hollow Knight.} Qui la discussione con alcuni
studenti devia su \emph{Hollow Knight} e \emph{Silksong}: il docente non ha
completato una delle prove più dure (l'ha iniziata e ha detto basta), e non
capisce il senso di un premio che è solo ``la medaglia d'onore'' di averla
fatta; uno studente difende la parte narrativa, legata ai rapporti tra padre
e figlio, che il docente trova sepolta dietro contenuti opzionali, legati
peraltro al finale migliore e non a quello normale. È anche il motivo per cui
ha abbandonato \emph{Hollow Knight} un paio di volte prima di finirlo, oltre a
una fase di \emph{backtracking} poco chiara; lui che, assicura, non si è mai
bloccato in un metroidvania in vita sua. ``Piccolo odio'', anzi, ``odio, ci
mancherebbe, avercene''.\footnote{Il passaggio è molto frammentato nella
trascrizione, con più voci sovrapposte; se ne riporta solo il senso
generale.}

\paragraph{Tutte le abilità dall'inizio.} Lo skill based è legato al
miglioramento delle skill del giocatore: è l'efficienza di cui si parlava. Si
possono anche dare al giocatore tutte le abilità fin dall'inizio, come in
\emph{Street Fighter} o in \emph{Smash Bros.}: l'obiettivo del giocatore sta
nello sviluppare \emph{mastery} nel sistema.

\paragraph{Easy fun.} Dall'altra parte c'è l'\emph{easy fun}, che nella
definizione di Lazzaro è \textbf{curiosità, novità, immaginazione}; anche
l'abilità intuitiva di risolvere cose, tendenzialmente legata ai puzzle, ma
non solo. Ci sta \emph{Oxenfree}, avventura narrativa horror, un \emph{teen
drama}: tutte le decisioni di gameplay, anche nei dialoghi, sono molto
intuitive. Vuoi fare lo stronzo con tuo fratello o vuoi fare \emph{bonding}
con lui? Vuoi aiutare i tuoi amici a diventare una coppia o li vuoi separare?
Ci stanno anche \emph{Animal Crossing} e \emph{Stardew Valley}, giochi con una
\textbf{bassa barriera di skill}; anche se \emph{Stardew Valley} forse non sta
proprio lì, perché c'è una quantità di informazioni così elevata, e spesso non
spiegata, che un certo tipo di ricerca, non solo di intuizione, è richiesto.
``Comunque, ripeto, è uno spettro.''

\paragraph{Approfondimento: le quattro chiavi di Lazzaro.} Bateman e Boon, in
\emph{21st Century Game Design}, usano lo studio di Nicole Lazzaro
(\emph{Why We Play Games: Four Keys to Emotion without Story}, 2004) per
interpretare i loro tipi di giocatore. L'\emph{hard fun} è il gameplay legato
a sfide, strategie e puzzle significativi: Lazzaro vi associa rabbia,
frustrazione e noia, ma soprattutto il \emph{fiero}, l'euforia del trionfo
personale su una difficoltà, ``quella sensazione che fa alzare le braccia al
cielo''. L'\emph{easy fun} è associato a meraviglia, stupore e mistero. Il
\emph{people factor} (il \emph{people fun}) è il divertimento sociale, in cui
la chiave è il piacere di stare con gli altri, schadenfreude compresa. La
quarta chiave, che il libro chiama ancora \emph{altered states} (gli stati
emotivi generati dalla percezione e dal comportamento), è quella che in
seguito Lazzaro ha rinominato \emph{serious fun}. È lo stesso vocabolario che
il modello di Quantic Foundry recupera per il polo basso della challenge.

\subsection{Fantasy}

La \textbf{fantasy} è la sospensione dell'incredulità. Da una parte la
\textbf{lore profonda}: un mondo molto profondo, tanta narrativa di base.
\emph{Mass Effect}, \emph{Dark Souls}, \emph{Hollow Knight}, \emph{Fallout},
\emph{Dragon Age}: giochi con una componente narrativa e una lore non
necessariamente alle spalle, ma anche nel mondo di gioco, quindi non solo
legata a un passato ma anche a un presente, una narrativa presente e molto
dominante.

Dall'altra parte il \textbf{setting generico o astratto}, o comunque giochi in
cui la narrativa è un \emph{flavor}, un \emph{dressing}, non una componente
importante del gameplay: non lo cambia. \emph{Street Fighter}, \emph{Gato
Roboto} (citato prima), \emph{Candy Crush}, \emph{World of Tanks} (ci sono i
carri armati che si sfidano, e basta); e anche un \emph{Call of Duty}, per il
tipo di motivazione che l'utente cerca. Al docente dà ragione il fatto che lì
dentro il single player ormai quasi non esista più: è più spinto in questa
direzione, soprattutto nel tempo, in una parabola discendente.

\subsection{Story}

Se la fantasy è la creazione del contorno del mondo di gioco, del contesto in
cui il giocatore gioca, la \textbf{story} è il \textbf{dramma}: esiste una
narrativa importante con un inizio e una fine nel mondo di gioco. Vale anche
per la profondità dei personaggi: un cast molto ben scritto, personaggi a
tutto tondo. Qui stanno tutti i giochi Telltale, i giochi BioWare (``riposi in
pace''), i \emph{Persona}, e in generale molti RPG con trame molto lunghe e
personaggi molto convoluti.

Dall'altra parte l'\textbf{open ended}: a volte la storia nel mondo di gioco
non è nemmeno spiegata, a volte è interpretata o ha un finale aperto.
\emph{Factorio} e \emph{SimCity} sono chiaramente open ended. \emph{Mario}
starebbe più verso il dramma, perché uno script esiste: liberi la principessa
da Bowser. Si può discutere sul fatto che non sia estremamente profondo né
abbia un grande cast di personaggi dalle molte personalità, ma ``è più qui
che qui''.

\subsection{Discovery}

La \textbf{discovery} è l'esplorazione, ciò che si diceva di
\emph{Minecraft}: la possibilità di esplorare qualcosa di sconosciuto, che a
volte si traduce proprio nell'azione fisica di andare oltre la collina, per
vedere che cosa c'è dall'altra parte. \emph{The Elder Scrolls}, \emph{Fallout},
\emph{The Legend of Zelda}, ma soprattutto gli open world come
\emph{Minecraft}, \emph{Valheim}, \emph{Skyrim}; e \emph{Subnautica}, quello
sott'acqua, ``bellissimo'': è curiosità.\footnote{Nella trascrizione il
secondo titolo compare come ``Venom''; la lettura \emph{Valheim} è la più
plausibile nel contesto, ma non certa.}

Dall'altra parte i giochi estremamente \textbf{pratici}, tecnici, con regole
definite, scriptate, e minime variabili nell'esplorazione. \emph{FIFA},
\emph{Madden} e i giochi sportivi hanno una discovery quasi a zero; così lo
\emph{Scarabeo}, gli \emph{scacchi} e i giochi da tavolo riprodotti.

\subsection{Design}

\textbf{Design} è un termine molto generico; come risposta emotiva e
motivazione generale si traduce nella \textbf{self-expression}: la
possibilità di modificare il modo di giocare, il personaggio principale,
secondo le proprie esigenze. Il personaggio con i capelli corti o lunghi,
rosa o verdi, le unghie nere, rosse o gialle; grasso, magro, alto, basso, e
così via.

Dall'altra parte il \textbf{personaggio fisso}. \emph{Mario} è \emph{fixed}:
nemmeno \emph{Mario Odyssey} lo è del tutto, perché si possono cambiare i
costumi, ma il personaggio resta quello. \emph{Braid} è all'estremo finale,
e anche \emph{Spelunky}; \emph{Spelunky 2} è ``un tacchettino'' più in là,
perché si può scegliere tra una decina di personaggi. In \emph{Super Mario
Galaxy 2} si poteva usare anche Luigi (il docente non ne è sicuro): è a metà
tra i due. Restano comunque personaggi fissi: l'asset è quello.

\subsection{A che cosa serve}

``Non dimenticatevi questo spettro'': è molto importante per definire le
\textbf{motivazioni principali} di un gioco, per \textbf{categorizzare
l'utente} e per chiedere all'utente che cosa gli interessa. Naturalmente,
le domande vanno fatte in un modo che abbia senso nel gergo del giocatore;
ma intanto si ha già questo framework, che si vedrà ibridato più avanti.

\paragraph{Il sito e il survey.} Sul sito di Quantic Foundry (il docente
mostra anche la versione vecchia, ``guarda che bello questo'') ci sono
\textbf{report sull'audience} e perfino una \textbf{perceptual mapping} che
funziona su una serie di giochi: ``e così si fa''. Ci sono molti dati che in
aula non si sono avuti. C'è anche un \textbf{survey} che si può compilare, e
da cui si possono eventualmente prendere domande (il docente non sa se sia
ancora aperto). Non solo categorizza negli archetipi di Quantic Foundry (che si
rivedranno, insieme ad altri più largamente condivisi), ma propone anche,
rispetto ai giochi preferiti, altri giochi vicini, che individui con
motivazioni simili hanno indicato. ``Strumento molto utile.''

\paragraph{Dati chiusi.} Purtroppo, come per altri framework, sondaggi,
\emph{audience report} e archetipizzazioni che si vedranno, i dati sono chiusi
o sono stati chiusi. Ci sono eccezioni, ma non le tratterà il docente:
arriveranno al terzo anno. ``Questo è il mantra, e ci devo stare pure io.''

% ---------------------------------------------------------------------
\section{Correlare: pattern, coincidenze e insight}
\label{sec:correlazione}

\subsection{Chiedersi perché}

L'ultimo punto della ricerca è la \textbf{correlazione}. Una volta che si hanno
tutti i dati, ci si chiede \textbf{perché} è successo qualcosa:
\begin{itemize}
  \item perché c'è questo buco di mercato?
  \item perché questo gioco ha fallito, rispetto a ciò che si è ottenuto?
  \item perché a questa persona non piace? Perché \emph{Concord} non è
        piaciuto all'utenza in generale? Com'è successo? Che cosa ha
        significato?
  \item c'è un \textbf{pattern}? Ogni tot anni esce questo tipo di gioco,
        emerge qualcosa, succede qualcosa? Come lo si può riprodurre? Come lo
        si può semplificare?
\end{itemize}
Per rispondere ad alcune di queste domande si sono già visti diversi metodi;
il più semplice sono i \textbf{nove perché} (sezione~\ref{sec:client-pitch}).

\subsection{Il cervello che vede pattern}

Nella slide c'è ``una grande immagine'': il docente si dichiara grande fan
delle teorie del complotto, di qualsiasi teoria del complotto aliena e non
solo.\footnote{In questo punto la trascrizione automatica entra in loop su una
parola ripetuta; il contenuto si recupera dalle frasi successive.} Non sono
vere, non hanno una profondità reale; ma tutti abbiamo questo meccanismo nel
cervello. È un \textbf{bias cognitivo}: riconosciamo pattern, molto bene o
molto male (per esempio, nella statistica), perché in passato questo ci ha
salvato la vita.

Se il nostro cervello fosse adattato a calcolare tutte le possibilità per cui,
uscendo dalla stanza la mattina, si può morire (un'auto che ti investe, un
asteroide in testa, un cane che ti azzanna alla gola), non usciremmo mai più
di casa. Il cervello si è adattato a questo, ma anche a riconoscere pattern
utili: le cose commestibili, le cose esteticamente appaganti. La
\textbf{simmetria} è così importante per noi perché generalmente le cose
simmetriche erano cose buone, per il nostro antenato nella vecchia caverna.

\paragraph{Lincoln e Kennedy.} Molte correlazioni nascono da questa lettura
sbagliata dei pattern. L'esempio della slide sono le famose coincidenze tra
Lincoln e Kennedy: entrambi assassinati, con dettagli che sembrano
ripetersi, come la moglie che avrebbe avvisato di non uscire quella mattina.
Sono tutte coincidenze (``magari qualcuno vuole crederci, ma vi assicuro che
sono coincidenze''): bisognerebbe investigare a fondo ciascun punto. Perché
sono morti entrambi di venerdì? È una coincidenza relativa, perché molti
eventi pubblici si tengono il venerdì, il sabato e la domenica, per
massimizzare l'affluenza.

\paragraph{Smontare le proprie ipotesi.} Quindi, una volta che si hanno i
dati, o si è trovato un \emph{blue ocean} (sezione~\ref{sec:utenti-mercato}),
magari con la perceptual mapping (sezione~\ref{sec:perceptual-mapping}),
``non ci andate subito, per favore''. Bisogna domandarsi, chiedere,
approfondire, \textbf{cercare di smontare le proprie ipotesi}.

\subsection{Dall'insight all'empatia}

Correlando e sintetizzando la ricerca si capiscono i \emph{perché}; i perché
portano all'\textbf{insight} visto prima (sezione~\ref{sec:framework-triadica}),
rispondono ai dubbi e sviluppano \textbf{empatia} con il target audience. Se si
ha insight sui dati, se si ha empatia con il target audience, se ci si è
chiesti abbastanza volte perché, allora si possono \textbf{soddisfare le
necessità e le emozioni del target audience}. È l'obiettivo finale,
considerando tutti i limiti in cui si ``nuota'': gli stakeholder, il lato
business, la visione, il gameplay. È ciò che si dovrebbe raggiungere alla
fine, e lo si ritroverà quando si parlerà di visione, in una versione
``rimaneggiata'' della stessa slide (``una spolverata di olio di gomito'').

% ---------------------------------------------------------------------
\section{La fine della pre-produzione}
\label{sec:fine-preproduzione}

Una volta ottenuti i dati, la fine della pre-produzione si sancisce con una
serie di \textbf{risposte} (tabella~\ref{tab:domande-visione}).
\begin{description}
  \item[Che cosa vogliamo fare?] Che cosa vogliamo fare noi, come team? Che
        cosa ci porta avanti, qual è la nostra motivazione, che giochi
        vogliamo fare?
  \item[Che cosa possiamo fare?] Qui si va verso la visione più produttiva:
        che cosa si può fare realisticamente, con quali abilità, che si hanno
        o che si vogliono (e si possono) sviluppare. Per quanto al docente
        piaccia la pixel art, o fare UI, dire che da oggi a domani, o da oggi
        alla fine della produzione, diventerà un pixel artist o un 2D artist è
        improponibile. \textbf{Siate concreti; piuttosto, sparate basso.} Qui
        dentro ci sono le abilità del team e la \textbf{fattibilità}
        (\emph{feasibility}), legata allo \emph{scope} e alla scala del
        progetto: che cosa si può fare, dati questi tempi, questi limiti,
        queste risorse.
  \item[Che cosa sappiamo che funziona?] Lo dicono il target audience e i
        competitor: sono i \emph{must}.
  \item[Che cosa sappiamo che non funziona?] Lo dicono il target audience e i
        \emph{must not}. Li si vedrà.
  \item[Che cosa vogliono gli stakeholder?] Lato manageriale e lato cliente
        (sezione~\ref{sec:stakeholder}).
  \item[Qual è il nostro business?] Rispetto al target audience, a livello di
        vendite: un gioco con le pubblicità oppure no? Si aggiungono DLC? Si
        vuole un gioco transmediale, altri prodotti derivati dall'IP?
        ``Dipende.''
\end{description}

\begin{table}[htbp]
  \centering
  \caption{Le domande che chiudono la pre-produzione e da cui nasce la game
  vision.}
  \label{tab:domande-visione}
  \begin{tabular}{@{}p{0.34\textwidth}p{0.58\textwidth}@{}}
    \toprule
    \textbf{Domanda} & \textbf{Da dove viene la risposta} \\
    \midrule
    Che cosa vogliamo fare? & motivazione e desideri del team \\
    Che cosa possiamo fare? & abilità del team, fattibilità, scope, tempi, risorse \\
    Che cosa sappiamo che funziona? & target audience e competitor (\emph{must}) \\
    Che cosa sappiamo che non funziona? & target audience (\emph{must not}) \\
    Che cosa vogliono gli stakeholder? & management e cliente \\
    Qual è il nostro business? & vendite: pubblicità, DLC, transmedialità \\
    \bottomrule
  \end{tabular}
\end{table}

\paragraph{La game vision.} L'idea che interconnette tutte queste risposte è la
\textbf{game vision}: in concreto, il documento di visione finale, più i
pilastri, più una serie di altre cose che si vedranno
(sezioni~\ref{sec:utenti-mercato} e~\ref{sec:principi-iterazione}). ``Questo
era per quanto riguarda la ricerca.''

% ---------------------------------------------------------------------
\section{Idee e brainstorming: i principi}
\label{sec:principi-brainstorming}

\subsection{Le idee sono inutili}

Si passa, ``veloci'', alle idee e al brainstorming. Anche qui arriva una serie
di strumenti, da non usare sempre tutti (``è una lista, e ce ne sono molti
altri''), ma che possono dare spunti su come discutere, con chi farlo e in che
tempi.

Il punto di partenza è una provocazione: \textbf{le idee sono inutili}. Non
esiste un'idea valida, che abbia senso nella produzione, se non è
\textbf{contestualizzata} nei limiti, nella visione, in tutto ciò che non si
può fare. Iterazione e implementazione sono tutto, in un progetto
(sezione~\ref{sec:principi-iterazione}). Le slide sono provocazioni, avverte
il docente: non dicono che cosa si deve fare esattamente, ma che cosa si
potrebbe fare.

\subsection{Generare tante idee, scartare le prime}

L'obiettivo di un brainstorming è \textbf{generare un ``macello'' di idee}: la
slide ne indica cento, ed è davvero tanto. Poi:
\begin{itemize}
  \item si scartano quelle troppo vicine all'\textbf{affezione del team}, a cui
        ci si affeziona troppo e che quindi non si riesce a valutare in modo
        obiettivo; anche questa è una provocazione, ma tendenzialmente si
        possono \textbf{scartare le prime idee} che escono, le prime venti;
  \item si valutano le altre con i dati che si hanno, i pro e i contro, e così
        via.
\end{itemize}

\paragraph{La legge del martello.} Perché scartare le prime venti? Andare
contro i bias cognitivi è sempre difficile, ma esiste un principio che,
semplificando molto, dice: \textbf{se hai un martello, vedrai tutti i problemi
come chiodi} (il termine ``chiodo'' arriva dalla classe). Siamo abituati a
risolvere i problemi con gli strumenti che usiamo di solito. ``Squadra che
vince non si cambia'': se qualcosa funziona, se vediamo il mondo in un certo
modo perché ci ha dato valore per anni o perché con quello risolviamo quasi
tutti i problemi che ci arrivano addosso, potremmo avere i paraocchi e non
vedere altro.

\paragraph{Non da soli.} \textbf{Lavorare in team} aiuta a superare questo
limite: io vedo il gioco, l'idea, il modo di risolvere i problemi in una certa
maniera, e altri lo vedono in un'altra. Il docente ha sempre trovato quasi
sufficiente lavorare con altri designer, con altre persone: fare
brainstorming \textbf{non da soli}, perché da soli è difficile. Chi vuole
farlo da solo, o fa fatica a mettersi nei panni degli altri o a ``lasciar
andare nel tritacarne'' le proprie idee, può trovare utili altri metodi, come
l'\emph{acting brainstorming} o le altre tecniche che si vedranno.

\paragraph{Approfondimento: prima da soli, poi insieme.} Schell, nei suoi
sedici consigli pratici per il brainstorming (\emph{The Art of Game Design},
capitolo~7), dedica l'ultimo al trovare un partner: in due si può arrivare a
grandi soluzioni molto più in fretta, rimbalzandosi le idee; ma aggiungere
persone non aiuta necessariamente, e di solito i gruppi migliori non superano
le quattro persone. I gruppi funzionano meglio su un problema
\textbf{circoscritto} che su uno ampio e aperto; e la ricerca, nota Schell,
mostra che un gruppo che si riunisce ``a freddo'' per fare brainstorming
spreca tempo: è molto meglio che ciascuno faccia prima brainstorming
\textbf{da solo}, e che solo dopo ci si riunisca per condividere, mescolare le
idee e risolvere problemi insieme. Vanno evitati come partner quelli che
bucano ogni idea o che hanno gusti molto ristretti.

\subsection{Le regole del brainstorming}

\paragraph{Comprendere la mente umana.} Per capire che cosa \emph{non} fare,
bisogna capire la mente umana: gli eventuali bias cognitivi, da mettere da
parte se ci sono.

\paragraph{Tutti sono creativi.} Bisogna trovare modi per far emergere il più
possibile tante idee. Generare cento idee in poco tempo non è semplice:
bisogna trovare il modo di ``spararne fuori'' il più possibile. Per chi non è
particolarmente creativo, le slide indicano due modi molto utili, e se ne
vedranno altri (sezione~\ref{sec:tecniche-idee}).

\paragraph{Lavorare dentro il framework, ma dopo.} Vale in due sensi. Quando,
un passo dopo il brainstorming, si comincia a setacciare le idee, vanno tolte
tutte quelle non attuabili, non realistiche da applicare. Ma \textbf{non lo si
fa prima}: durante il brainstorming non ci si frena, \textbf{anything goes},
si butta fuori qualsiasi cosa. Anche da qualcosa che poi verrà scartato, o che
non ha senso nel contesto della produzione, può emergere qualcosa di
interessante o unico. Dopo, quando si valuta, si ricordano la visione, le
limitazioni, i documenti, i limiti.

\paragraph{Rispettare tempi e ruoli.} Se il brainstorming ha regole di tempo o
di ruoli, le si rispetta: se sta parlando il docente non parla Alessandro, e
se sta parlando Lorenzo non parla il docente. L'obiettivo del brainstorming è
\textbf{dare a tutti la propria voce}, fare in modo che tutti esprimano la
propria idea. E, come sempre, se dopo (non in quel momento, molto dopo) la
propria idea non viene accettata, non significa essere un cattivo designer o
aver fatto un brutto lavoro: l'idea, quando viene messa a terra, passa in
mano a tutti. Durante il brainstorming tutti devono poter parlare, sentirsi
liberi di dire la propria e \textbf{non sentirsi giudicati}. Ci sarà il momento
di valutare attentamente, ragionare, setacciare, tagliare, ``bruciare'' idee,
ma non è quello.

\paragraph{Soluzioni \emph{future proof}.} Dopo, prima della produzione, durante
la pre-produzione, è importante essere \textbf{severi} per primi con ciò che
si idea e si documenta. Se non lo si fa lì, arriverà qualcuno, volutamente o
no, con una domanda o un dubbio che distrugge quel design, o che mostra una
situazione a cui non si era pensato, e quello è un problema. Bisogna avere
il design \textbf{più solido possibile all'inizio della produzione}: tenere le
soluzioni \emph{future proof} e scartare le altre. Se qualcosa convince fino
a un certo punto, ma ha dei limiti o delle pericolosità future, si valuta se
il rischio vale la candela; se no, via.

% ---------------------------------------------------------------------
\section{Il design elegante}
\label{sec:design-elegante}

\subsection{Che cosa vuol dire}

``Puntate al design elegante. Questo non è vero sempre, ragazzi.'' Fare design
elegante piace a tutti (``piacerà a voi quando saprete che cosa vuol dire''), ma
è difficile. Un design elegante è non solo \textbf{easy to learn}, ma
\textbf{difficult to master}; è \textbf{molto interconnesso}, con regole che
hanno ricadute su altre regole. Un tiro di dado che gestisce sia la vita sia
la stamina, una statistica che gestisce due cose diverse, un attacco che fa
due cose: è elegante, perché è molto \textbf{condensato}.

\paragraph{Il rovescio della medaglia.} Proprio per questo, se in futuro si
deve toccare quella statistica, la vita o il sistema di dadi, l'intervento
impatta su entrambe le cose. Se non lo si ``blinda'', se non ci si pensa
molto bene all'inizio, ``vi ritorna indietro due volte sui denti'': fa male il
doppio. È il legame con le soluzioni \emph{future proof}.

\subsection{Elegante non vuol dire migliore}

Non c'è sempre bisogno di fare design elegante. Ottimi giochi, soprattutto
ottimi giochi da tavolo, lo fanno: \emph{Root}, per esempio. Ma un ottimo gioco
da tavolo che non è minimamente elegante è \emph{Dune}, nella riedizione di
Gale Force Nine: uno dei giochi da tavolo preferiti del docente, ``non
minimamente elegante di design, però funziona estremamente bene''; è anche
vero che è figlio di un'altra epoca, di un altro grande gioco.\footnote{La
riedizione di Gale Force Nine (2019) riprende il \emph{Dune} di Avalon Hill
del 1979. La frase nella trascrizione è disturbata (``figlio di un'altra equa,
grande gioco'').}

\paragraph{Gli esempi della classe.} Gli studenti propongono altri titoli:
\begin{itemize}
  \item \emph{Here to Slay}, che uno studente trova elegante per la sua
        condizione di vittoria legata ai mostri;
  \item \emph{Unstable Unicorns}, che al docente era venuto in mente prima
        perché è ancora più elegante, anche se con meno meccaniche: ``è
        \emph{Yu-Gi-Oh!} semplificato, \emph{Yu-Gi-Oh!} per \emph{toddlers}'';
  \item \emph{7 Wonders}, con un sistema di \emph{drafting} molto interessante,
        anche se il docente non lo considera elegantissimo (forse,
        considerando le espansioni, qualcosa si può pensare);
  \item il drafting di \emph{Terraforming Mars} (confuso per un attimo con
        \emph{Surviving Mars}: ``uno è il gioco da tavolo e l'altro è il gioco
        digitale'').
\end{itemize}
Anche i giochi molto contenuti, con scatole ridotte, possono essere
estremamente eleganti; al docente sfugge il nome di un gioco ``con
l'animale'', ``perduto per sempre; tanto mi ritornerebbe prima o poi''.

\paragraph{Dalle meccaniche ai sistemi.} Uno studente collega il discorso al
design di sistema. Il docente conferma: qui si parla di meccaniche, ma nella
massima estensione i sistemi interconnessi diventano il \textbf{system
design}, se il gioco è sistemico (e non è detto che lo sia). In generale, il
design elegante è ``una buona \emph{forma mentis}''.

\paragraph{Approfondimento: la lente dell'eleganza.} Schell dà un modo pratico
per misurare l'eleganza: \textbf{contare gli scopi} di ciascun elemento. I
puntini di \emph{Pac-Man} danno al giocatore un obiettivo a breve termine
(mangiare quelli vicini) e uno a lungo termine (ripulire il livello); lo
rallentano leggermente quando li mangia, rendendo più rischiosi i corridoi
pieni; danno punti, misura del successo; e i punti fanno guadagnare vite
extra. Cinque scopi per un elemento semplicissimo. La \emph{Lens of Elegance}
(lente n.~49) chiede: quali sono gli elementi del mio gioco? Quali scopi ha
ciascuno? Quelli con uno o due scopi si possono combinare tra loro o
eliminare? Quelli con molti scopi possono assumerne altri ancora? Subito
dopo, però, Schell introduce la \emph{Lens of Character} (lente n.~50):
eleganza e carattere sono opposti da tenere in equilibrio. La torre di Pisa,
resa ``elegante'' raddrizzandola, non la visiterebbe nessuno; le pedine di
\emph{Monopoly} non c'entrano nulla con il mercato immobiliare, e Mario non ha
alcun motivo di essere un idraulico, ma è proprio questo che dà loro
carattere. È il senso del ``non sempre'' del docente.

% ---------------------------------------------------------------------
\section{Essere creativi}
\label{sec:creativita}

Come essere creativi? Non c'è una vera soluzione, e non è importante se si è
più o meno portati per la creatività. Conviene però identificare eventuali
\textbf{momenti o rituali} in cui si pensa meglio.

\paragraph{I rituali.} Woody Allen, racconta il docente, non riusciva a
scrivere le sceneggiature dei suoi film se nel suo albergo non c'era una
doccia: doveva per forza averla. Un altro esempio citato è quello di chi
faceva un sonnellino e, al risveglio, aveva l'idea.\footnote{Il secondo nome è
storpiato nella trascrizione (``So, Cristiano''); l'aneddoto del sonnellino
non è attribuibile con certezza.} Ci possono essere momenti di calma, o al
contrario di estrema agitazione, in cui si è più prolifici.

\paragraph{La praticità.} Il docente si trova molto più comodo a scrivere,
come si vede, alla lavagna o sui fogli di carta, a scarabocchiare: la
praticità lo aiuta, ad altri potrebbe dare fastidio. Racconta di un altro
senior con cui lavorava, anche lui molto pratico ma con i fogli digitali:
quando si parlava di combat, lui disegnava gli \emph{stickman} che
risolvevano il problema, mentre il docente era davanti alla camera ``a
prendere a bastonate il nulla''. Non è obbligatorio trovare un metodo; ma se
si ha bisogno di stimolare la creatività, conviene identificare i momenti in
cui si è più prolifici.

% ---------------------------------------------------------------------
\section{Le tecniche per generare idee}
\label{sec:tecniche-idee}

\subsection{Una rassegna, non un obbligo}

Il docente procede ``a ruota'', fino a un certo punto, rimandando il resto al
venerdì. Non si sofferma molto sui singoli punti, perché sono tutti modi di
mettere giù un'idea, fare brainstorming o creare qualcosa, e non vanno
applicati tutti: possono essere utili uno adesso, un altro tra vent'anni.
Come sempre, sono strumenti, e ce n'è ``davvero una miriade''. In un
brainstorming non codificato ognuno può fare quello che vuole; se è
codificato, ci saranno limiti, tempistiche e un \emph{handler} che decide chi
parla e quando, ``non perché è brutto e cattivo'', ma per garantire che
parlino tutti.

\subsection{Le categorie sensoriali}

Il primo modo divide la questione in \textbf{categorie}: teoricamente tutto si
può scomporre in categoria \textbf{visiva}, \textbf{uditiva}, \textbf{tattile},
\textbf{cinestetica}, \textbf{spaziale}, \textbf{enattiva} (la messa in moto,
l'azione effettiva), \textbf{iconica} (popolare, conosciuta) e
\textbf{simbolica}. L'esempio della slide è la musica:
\begin{description}
  \item[Visiva] tutto ciò che è visualizzazione: il DJ che mostra le curve dei
        bassi e degli alti.
  \item[Simbolica] uno spartito, la chiave di violino, le note: tutta la
        simbologia legata specificamente alla musica, che è anche visiva
        quando la si mette per iscritto.
  \item[Uditiva] ``manco lo dico''.
  \item[Tattile] sentire i bassi: le persone sorde sentono i bassi delle
        canzoni.
  \item[Cinestetica] ballare, attuare quell'azione.
  \item[Spaziale] lo spazio occupato da un'orchestra, o da una band rispetto al
        pubblico.
  \item[Enattiva] suonare lo strumento.
  \item[Iconica] un cantante pop, un'immagine musicale molto nota.
\end{description}
Si dovrebbe poterlo fare con qualsiasi cosa.

\subsection{Opposite direction}

È una delle tecniche preferite del docente.
\begin{enumerate}
  \item Si prende un gioco o un'idea che piace. Se si è già oltre la ricerca,
        si ha già una lista di categorie, generi e competitor possibili.
  \item Se ne scrivono le componenti o le \emph{game qualities} principali.
  \item Si trova il \textbf{gioco opposto}.
  \item Si uniscono i due: non gli elementi, ma i \textbf{due giochi
        insieme}.
\end{enumerate}
L'esempio: \emph{Doom} è FPS, \emph{fast paced}, \emph{gory}; l'opposto è
terza persona, \emph{slow paced}, \emph{cute}: \emph{Animal Crossing}. Insieme:
un gioco con un'arte simile ad \emph{Animal Crossing}, ma con una palette meno
satura, in cui di giorno si costruisce il proprio villaggio, con cagnolini
``carini e pucciosi'' come personaggi, e di notte il villaggio viene assalito
da orde di bestie immonde a cui si spara con un fucile a canne mozze (``sono
uscito un attimo dal seminato'').

\paragraph{Si venderebbe.} Probabilmente si potrebbe vendere, anche perché
quando \emph{Doom} uscì insieme ad \emph{Animal Crossing} la community mise
insieme Isabelle e il Doom Slayer.\footnote{Il docente esita tra \emph{Doom}
(2016) e il suo seguito: si tratta di \emph{Doom Eternal}, uscito lo stesso
giorno di \emph{Animal Crossing: New Horizons} (20 marzo 2020).} Ci sono un
sacco di fan art e animazioni con Isabelle che gira con la doppietta a
uccidere mostri e il Doom Slayer che cura il suo giardinetto e annaffia le
piante tranquillamente. La tecnica funziona soprattutto per creare cose
\textbf{innovative} o potenzialmente mai viste: \emph{Palworld}, secondo il
docente, è stato creato sicuramente così (si veda anche l'opposizione,
sezione~\ref{sec:competitor-diretti}).

\subsection{Break the theme}

Con la stessa idea di fondo, \textbf{rompere il tema}: si prende un'idea, se ne
elencano le \textbf{tre componenti dominanti} e se ne fa il contrario.
L'esempio: un ristorante serve il cibo, lo si paga, e lo chef cucina per te.
Un ristorante che non serve cibo, in cui non si paga il cibo e lo chef non
cucina per te, probabilmente non sarebbe un ristorante; ma togliendo alcune di
queste componenti si arriva a luoghi in cui si portano gli alimenti e si
cucina da sé, con eventualmente degli aiutanti. Il docente non mangerebbe mai
così, ma il punto è un altro: non è necessario prendere tutti i punti, se ne
possono prendere solo alcuni, e trovare così un potenziale nuovo gioco.

\paragraph{Approfondimento: distruggere le proprie assunzioni.} Tra i consigli
per il brainstorming di Schell ce n'è uno molto vicino, che lui attribuisce
al designer Rob Daviau: fare una lista di tutto ciò che si \emph{assume} vero
sul proprio gioco (``assumo che si giochi al chiuso'', ``assumo che il
giocatore guardi lo schermo'', ``assumo che usi un solo dito sul
touchscreen''), e poi chiedersi, voce per voce, come funzionerebbe il gioco se
quell'assunzione non fosse vera. Quasi sempre l'assunzione resta; ma ogni
tanto farla saltare dà un'intuizione decisiva. Daviau è arrivato a \emph{Risk
Legacy} distruggendo l'assunzione standard dei giochi da tavolo: ``una partita
non influenza la successiva''.

\subsection{Le regole euristiche}

Le \textbf{regole euristiche} sono una serie di regole che fanno leva su
alcuni bias cognitivi; nel complesso sono ``veramente complesse'', e ce ne
sono molte altre oltre a quelle della slide. Alcune delle principali:
\begin{itemize}
  \item se hai difficoltà a identificare un problema, \textbf{visualizzalo}:
        scrivilo, scarabocchialo;
  \item se non hai una soluzione, prova ad \textbf{assumere una soluzione} e
        parti da lì (il cosiddetto \emph{bottom-up design});
  \item se hai un problema concreto, \textbf{astrai}; se hai un problema
        astratto, \textbf{concretizza}.
\end{itemize}
Sono questioni già largamente note.

\subsection{What if}

Si scrivono domande del tipo \emph{what if}: ``e se\dots?''. Le slide ne
riportano una lista. Il docente spiega perché la trova ``una figata'' con i
fumetti: non è un assiduo lettore di Marvel o DC, ma divora qualsiasi cosa a
fumetti. \emph{Superman: Red Son} è il \emph{what if} di Superman se
atterrasse nell'Unione Sovietica di Stalin: diventa un compagno di Stalin, un
superuomo stalinista, e l'America, senza il suo protettore, deve fare la
corsa agli armamenti con Lex Luthor.\footnote{Nella trascrizione il titolo
compare come ``Superman 9 Son''.} Molte delle storie più amate di entrambe le
case sono \emph{what if}: la Marvel li chiama proprio \emph{What If}, la DC
\emph{Elseworlds}. C'è anche il \emph{what if} Marvel con gli zombie, scritto
da Robert Kirkman, da cui è stata tratta una serie.\footnote{\emph{Marvel
Zombies}; nella trascrizione l'autore compare come ``Kierkegaard''.} ``Si può
creare molto così.''

\subsection{Wicked questions}

Molto simili ai \emph{what if}: ci si chiede \textbf{come può la soluzione al
nostro problema essere X e Y contemporaneamente}, dove X e Y sono soluzioni
proposte dal gruppo e scritte su foglietti, messi in un grande calderone da
cui si pesca.\footnote{Le \emph{wicked questions} sono una delle
\emph{Liberating Structures} (sezione~\ref{sec:problema-necessita}).}

\subsection{SIT e TRIZ}

Il \textbf{SIT} (\emph{Systematic Inventive Thinking}) è simile all'euristica e
ha le sue regole, elencate nelle slide:
\begin{itemize}
  \item \textbf{sottrazione}, \textbf{moltiplicazione}, \textbf{divisione}:
        prendere un problema e scomporlo nelle sue componenti principali;
  \item \textbf{unificazione}: è così che funziona il design elegante
        (sezione~\ref{sec:design-elegante});
  \item il \textbf{closed world}: il concetto per cui, quando si formula il
        problema, la soluzione risiede nel problema stesso; è solo una
        questione di formulazione, o di vederlo da un altro punto di vista.
\end{itemize}
Il concetto di \emph{closed world} è alla base di un altro framework,
estremamente complesso, che si vedrà: il \textbf{TRIZ}.\footnote{Nella
trascrizione compare come ``Trits''. Il TRIZ è la teoria per la soluzione
inventiva dei problemi sviluppata da Genrich Altshuller; il SIT ne è una
derivazione semplificata.}

\subsection{Limitation}

Si prende il proprio gioco, o un genere, e gli si impongono \textbf{tre forti
limitazioni}. Un soulslike senza combattimento melee, senza schivata e senza
atmosfera cupa: sarebbe ancora un soulslike? Probabilmente sì.
\emph{Enotria}, per esempio, ha fatto una di queste tre limitazioni: niente
\emph{dark mood}. Anche le build permanenti, nota il docente, non sono
necessariamente uno dei punti identificativi di \emph{Bloodborne} o di
\emph{Dark Souls}.\footnote{Nella trascrizione il primo titolo compare come
``Blackboard''; la lettura \emph{Bloodborne} è la più plausibile.}

\subsection{Drawing e storyboard}

Si disegna: \emph{stickman}, nessuna parola. A volte si possono usare solo
figure con un determinato simbolo, come il cerchio. Una variante è lo
\textbf{storyboard}: si prende un foglio già diviso nei riquadri dello
storyboard e ci si disegna dentro una storiella. Può servire a trovare
qualcosa di non visto.

\subsection{Sharing}

Si parla con le persone, si condividono le cose. Nella slide c'è una mamma che
ha fatto una maglia incredibile: ci sono persone che fanno cose che non si
direbbe mai, o cose molto lontane dalla propria sensibilità. Più si conosce,
meglio è, anche perché \textbf{si può fare un gioco su qualsiasi cosa},
veramente qualsiasi. Bisogna essere \textbf{curiosi} e fare domande.

Spesso è anche un obbligo: se si viene assunti per fare un gioco ambientato
nell'antica Roma, bisogna conoscerla; verrà dato del materiale, ma bisognerà
anche fare ricerca per conto proprio su usi e costumi. Non si sfugge. Ma si
potrebbe anche voler partire da concetti lontani dalla sensibilità comune, o
inserire cose curiose che possono catturare un certo tipo di utenza.

\paragraph{Approfondimento: lo \emph{slow boil}.} In \emph{Challenges for Game
Designers}, Brathwaite e Schreiber descrivono diversi approcci al game design.
Uno è il \emph{blue sky}, senza altri vincoli che l'immaginazione del
designer (raro, perché di solito si lavora sotto vincoli esterni). Un altro è
quello che chiamano \emph{slow boil}: dato un tema e un'ambientazione (il loro
esempio è ``gli irlandesi nella Boston del 1880''), il designer avvia una
massiccia missione di ricerca, con poca direzione purché copra il tema. Come
dice uno di loro, ``mi riempio la testa di tutto quello che trovo
sull'argomento, e poi aspetto: un design emergerà''. Reiner Knizia, ricordano,
ha regolarmente una trentina di progetti in corso e preferisce rigirarli a
lungo nella mente prima di passare al prototipo; e un errore comune dei
designer alle prime armi è forzare un gioco troppo presto. È la stessa
curiosità che il docente raccomanda con lo \emph{sharing}.

\subsection{Ad absurdum}

Si prendono il problema o la necessità, il contesto, gli stakeholder e così
via, e li si \textbf{spinge al massimo} e \textbf{al minimo}; si risolve il
caso estremo nei due sensi. Teoricamente, poi, si valuta la via di mezzo (si
veda anche il ricorso alle situazioni estreme nella client pitch direction
analysis, sezione~\ref{sec:client-pitch}).

\subsection{Che cosa resta da vedere}

Resta il \emph{come} si fa brainstorming, ma il docente decide di lasciarlo
al venerdì: ``abbiamo quasi concluso''. A quel punto dà il task, ``perché
altrimenti mi restate a bocca asciutta per troppo tempo''.

% ---------------------------------------------------------------------
\section{Il task della settimana}
\label{sec:task-gd5}

Il brainstorming non è ancora stato visto per intero, ma la slide c'è, e se ne
riparlerà. Il task si fa \textbf{a gruppi}: gruppi di due, tranne uno di tre.

\begin{description}
  \item[I gruppi] Francesca e Antonin\footnote{Il nome compare così nella
        trascrizione; potrebbe essere troncato (Antonino, Antonio).};
        Federica e Alberto; Francesco e Lorenzo; Valerio, Mattia e Daniele.
        Il docente si confonde un paio di volte nell'assegnarli (``quattro
        gruppi, oggi veramente fuori di testa'') e chiede che qualcuno gli
        rimandi i gruppi per messaggio.
  \item[I giochi] I primi due gruppi lavorano su \textbf{Marvel Snap}; gli
        ultimi due, compreso quello da tre (``vi beccate la roba più
        difficile''), su \textbf{Yu-Gi-Oh! Master Duel}. Sono due
        \textbf{digital card game}, non TCG, perché le carte non si possono
        scambiare; \emph{Marvel Snap} è il più facile. Per \emph{Master Duel}
        c'è una sola banlist, ``credo''. La ricerca va naturalmente limitata a
        questi giochi: è già un grosso lavoro.
  \item[La ricerca] Tramite ricerca su internet (Reddit, Discord, dove si
        vuole) bisogna identificare \textbf{almeno tre problemi o necessità}:
        tre \emph{pain point} che il target audience di riferimento ha
        individuato come problematici. Vanno trovati i dati (il feedback),
        analizzati, spacchettati: qual è il problema?
  \item[L'analisi] Il problema va analizzato il più possibile con gli
        strumenti visti: identificato ed eventualmente categorizzato (con
        Cynefin, tabella~\ref{tab:cynefin}, o con le Liberating Structures,
        sezione~\ref{sec:problema-necessita}). Si può usare la perceptual
        mapping (sezione~\ref{sec:perceptual-mapping}), anche per vedere come
        i competitor vicini hanno eventualmente risolto la stessa cosa. In
        ogni caso va documentata la parte di \textbf{dati}, la prima parte
        della lezione.
  \item[La soluzione] Poi si genera una soluzione. Il consiglio del docente è
        usare una delle tecniche più semplici: \textbf{opposite direction},
        \textbf{limitation} o \textbf{break the theme}
        (sezione~\ref{sec:tecniche-idee}). Sconsiglia SIT ed euristica,
        perché richiederebbero tanto lavoro di ricerca, tempo che è meglio
        dedicare a capire che cosa dicono i dati dell'utenza.
  \item[Il documento di brainstorming] Ci deve essere un \textbf{documento di
        brainstorming} in cui si butta giù di tutto (``permettetemi il
        francese, cazzate''), ci si ``tira la pizza addosso''. Quel documento
        \textbf{non deve seguire le golden rules}
        (sezione~\ref{sec:golden-rules}) e resta agli studenti; ma va creato,
        per mostrare di aver lavorato. Il docente non lo leggerà
        prontamente.
  \item[Il documento consegnato] Il docente leggerà il documento prodotto,
        composto da: i \textbf{tre pain point} dell'utente, \textbf{come}
        sono stati trovati, i \textbf{dati} spiegati e il \textbf{design} per
        risolvere ciascuno dei tre.
  \item[Scadenza] Il docente concede un giorno in più del solito:
        \textbf{venerdì della settimana successiva}, due giorni in più, anche
        perché nella lezione sul brainstorming arriveranno altre informazioni
        da inserire. Intanto si può cominciare a fare ricerca, ragionare,
        scambiarsi i numeri di telefono e organizzarsi.
\end{description}

La lezione si chiude con le ultime domande e con l'appuntamento a venerdì,
\textbf{in presenza}, dove si concluderà la parte su come fare brainstorming;
sciopero permettendo, con la previsione che ci saranno comunque ritardi nei
treni.

% ---------------------------------------------------------------------
\section*{Punti chiave del capitolo}
\addcontentsline{toc}{section}{Punti chiave del capitolo}

\begin{itemize}
  \item Il \textbf{production dictionary} non ha limiti di pagine nelle
        produzioni vere; può essere un file a sé o un paragrafo d'apertura
        ibridato in un altro documento; i termini li sceglie il lead, li
        riempiono molte mani, e ce n'è \textbf{uno per progetto}.
  \item \textbf{Documenti di visione} (production dictionary, documento di
        visione, pitch) contro \textbf{documenti di produzione} (i GDD). Il
        GDD unico non esiste più: dalla ``Bibbia'' si è passati alle wiki e a
        strutture che smistano i lettori verso i documenti giusti, curate
        dalla produzione.
  \item Deming: il ciclo PDCA e le \textbf{sette malattie mortali}, tra cui
        dirigere un'azienda basandosi solo sui dati visibili.
  \item I framework di design (core elements, game qualities, Bura) servono
        anche ad \textbf{analisi} e \textbf{pitching}.
  \item L'\textbf{elicitazione triadica}: tre competitor, il target audience
        raggruppa i due più simili e spiega perché; si cerca l'elemento
        dominante, i pain point e l'insight.
  \item La \textbf{perceptual mapping}: assi come spettri, per trovare
        combinazioni innovative e buchi di mercato; un vuoto non è un invito;
        \textbf{un grafico non basta}. Utile anche nel design (nemici: DPS,
        HP, velocità, melee/ranged) e nel brainstorming.
  \item L'\textbf{analisi degli utenti} è tra i punti più difficili: dati
        quantitativi, gusti, routine; target dominante e demografiche
        attorno; il feedback online è \textbf{polarizzato}; nessun gioco è per
        tutti.
  \item Survey: brevi (\textbf{meno di cinque minuti}), domande semplici e
        puntuali, molto più di cento risposte; i paganti li finiscono. Poi
        archetipi, osservazione e la mappa \emph{dicono, fanno, pensano,
        provano}. \textbf{Senza empatia non c'è design}.
  \item L'\textbf{analisi di mercato}: spazio, dimensione, saturazione, stato
        emotivo (soulslike saturi, monster collectible scontenti), trend
        (sempre in ritardo), barriere, tasse, logistica.
  \item Le \textbf{dodici motivazioni} di Quantic Foundry in sei gruppi
        (action, social, mastery, achievement, immersion, creativity), ognuna
        uno spettro; si descrive un gioco con le \textbf{dominanti}. L'unico
        \emph{fun} ammesso è quello delle quattro chiavi di Lazzaro.
  \item \textbf{Correlare} i dati senza cadere nei falsi pattern; smontare le
        proprie ipotesi; arrivare a insight ed empatia.
  \item La pre-produzione si chiude con le domande su che cosa si vuole e si
        può fare, che cosa funziona e non funziona, che cosa vogliono gli
        stakeholder e il business: la risposta è la \textbf{game vision}.
  \item Le idee sono inutili senza contesto: generarne tante, scartare le
        prime, attenzione alla \textbf{legge del martello}, brainstorming in
        gruppo e senza giudizio, poi severità e soluzioni \emph{future
        proof}. Il \textbf{design elegante} è desiderabile ma non obbligatorio.
  \item Tecniche: categorie sensoriali, \textbf{opposite direction},
        \textbf{break the theme}, euristica, what if, wicked questions, SIT e
        TRIZ, \textbf{limitation}, drawing e storyboard, sharing, ad absurdum.
  \item Task a gruppi: tre pain point di \emph{Marvel Snap} o \emph{Master
        Duel}, con dati, analisi e soluzione (opposite direction, limitation o
        break the theme), più un documento di brainstorming; scadenza venerdì
        della settimana successiva.
\end{itemize}
